\documentclass[12pt, a4paper, oneside]{report}

% ,  PACKAGES , 
\usepackage[margin=2cm]{geometry}
\usepackage[utf8]{inputenc}
\usepackage{amsmath, amssymb, amsthm, bm}
\usepackage{graphicx}
\usepackage[colorlinks=true, linkcolor=blue, citecolor=blue, urlcolor=blue]{hyperref}
\usepackage{booktabs}
\usepackage{setspace}
\usepackage{caption}
\usepackage{subcaption}
\usepackage{algorithm}
\usepackage{algpseudocode}
\usepackage{physics}
\usepackage{qcircuit}
\usepackage{tcolorbox}
\usepackage{float}
\usepackage{cite}
\onehalfspacing

\begin{document}

% ==========================================
% TITLE PAGE
% ==========================================
%TC:ignore
\begin{titlepage}
    \centering
    \vspace*{1cm}
    \includegraphics[width=0.5\textwidth]{imperial_logo.png} \\[1.5cm]
    {\large \textbf{Project Code: SPC-Appelbe-1}} \\[0.5cm]
    {\huge \bfseries Q-HAS: Quantum-Hierarchical Adaptive Steering} \\[0.5cm]
    {\Large \bfseries A Predictive Hybrid Framework for Topological Instability Detection in Plasma} \\[2cm]

    \textbf{Student CID:} 06068471 \\[0.5cm]
    \textbf{Supervisor:} Dr.\ Brian Appelbe \\
    \textbf{Assessor:} Dr.\ Aidan Crilly \\[2cm]

    \textbf{Word Count}: 9 584 \\

    \vfill
    Academic Year 2025--2026
\end{titlepage}
%TC:endignore

% ==========================================
% LAYPERSON’S SUMMARY
% ==========================================
%TC:ignore
\chapter*{Layperson’s Summary}
\addcontentsline{toc}{chapter}{Layperson’s Summary}

\begin{center}
{\Large \textbf{Quantum algorithms: a new tool for predicting plasma instabilities?}}
\end{center}

\noindent Fusion energy, the same process that powers the Sun, appears to be a clean and extremely promising source of energy. To achieve it on Earth, scientists must confine plasma, a gas so hot that its atoms break apart into electrically charged particles, inside powerful magnetic fields. The problem is that plasma is often chaotic: it twists, flows and can suddenly become unstable. Small disturbances can develop into larger ones, which are capable of damaging the walls of a reactor. Predicting when and where these instabilities will occur is one of the main challenges in plasma physics.

To study plasma, scientists run large-scale computer simulations. The plasma volume is divided into millions of small cells, and the equations governing its motion are solved in every cell at each instant. Smaller cells yield more accurate results, but also demand far more computing power. Using the finest resolution everywhere would require months of supercomputer time for a single simulation.

A common solution is \textit{Adaptive Mesh Refinement} (AMR): the computer increases resolution only where complex behaviour is detected. However, current AMR methods are too reactive, they refine the grid only after an instability has already grown large enough to be noticed. By that point, it is possible that the simulation has already lost important detail.

This project asks a different question: rather than reacting to instability, can we anticipate it? The approach investigated here uses a small quantum circuit as an early-warning system. Quantum computers process information using quantum bits (qubits), which can represent multiple possibilities at once through a property called \textit{superposition}. Another property, \textit{interference}, allows qubits to amplify useful signals while cancelling out noise (like the way noise-cancelling headphones work). These two properties, used in the right way, can provide a huge advantage over classical methods.

In the Q-HAS (Quantum-Hierarchical Adaptive Steering) framework developed in this project, indicators of plasma behaviour, such as the rate of change of flow velocity and magnetic field, are encoded into a small set of qubits. A quantum algorithm then evaluates whether the pattern of encoded signals is consistent with the early stages of a known instability, such as a shearing flow, a forming vortex, or a magnetic reconnection event. The quantum circuit does not directly solve the plasma equations; instead, it only guides the classical simulation by flagging regions that are likely to become unstable in the near future.

The result is a hybrid architecture: a classical supercomputer performs the detailed plasma calculations, while the quantum component acts as a predictive steering mechanism that focuses computational effort where it is most needed. Both a theoretical analysis and a proof-of-concept software implementation are presented, demonstrating the framework on eight standard plasma test cases including the Kelvin-Helmholtz instability, the Orszag-Tang vortex, Harris current sheets, and isolated vortex, reconnection, and rotor scenarios.

Although today’s quantum hardware remains limited in scale, this work provides a foundation for integrating quantum sensing into classical simulations. Like many modern quantum applications, the aim is not to replace classical computing architectures but to enhance their performance using quantum algorithms.

%TC:endignore
% ==========================================
% ABSTRACT
% ==========================================
\begin{abstract}
This report presents a novel framework to perform Adaptive Mesh Refinement (AMR) in MHD simulations. Quantum-Hierarchical Adaptive Steering (Q-HAS for short) is a hybrid quantum-classical approach. Its main objective is instability detection by minimising a structured Hamiltonian describing plasma anomalies on a coarse-grained flux graph. Each qubit's polar angle $\theta$ encodes a classical multi-indicator score for the instability risk of an edge of the graph. A QAOA-based algorithm evaluates a cost function to locate the riskiest regions of the graph. The framework is implemented with a 4th-order finite-difference MHD solver, a Qiskit-based VQA pipeline, and Optuna Bayesian training. Evaluation on four scenarios ($N = 256$, 8 qubits) gives a 0.66\% composite-loss advantage of Q-HAS over classical AMR, statistically significant in favour of Q-HAS only on Kelvin--Helmholtz ($p = 0.026$) and Harris Tearing ($p = 0.001$); and statistically significant in favour of classical AMR on Orszag--Tang (OT) and MHD Rotor.  Topological attribution reveals $+3.8$pp and $+5.3$pp Q-HAS advantages in topology-rich sub-regions of Tearing and OT respectively. On aggregate per-cell decision accuracy, however, classical AMR is correct on 46.6\% of cells vs 43.3\% for Q-HAS, and achieves lower L2 error vs DNS on all four scenarios. Whether the localised advantage scales to larger VQA grids or higher Reynolds numbers remains an open question. All circuits run on the Qiskit-Aer simulator.
\end{abstract}

\tableofcontents

\chapter{Introduction}
\label{chap:introduction}

Magnetohydrodynamics (MHD)~\cite{freidberg2014,goedbloed2004} governs the nonlinear coupling of fluid motion and electromagnetism in electrically conducting plasmas. In regimes ranging from magnetic confinement fusion~\cite{wesson2011} to space plasmas and astrophysical flows, small perturbations evolve into coherent instabilities: vortices, current sheets, magnetic reconnection events, and kink modes. Accurately predicting and controlling these structures helps us improve our understanding of complex physics systems.

The MHD equations combine the Navier--Stokes equations with Maxwell's equations. In the incompressible MHD case, the magnetic field evolves through advection and resistive diffusion. Different types of anomalies arise from the nonlinear coupling between velocity and magnetic fields: shearing flows trigger the Kelvin--Helmholtz instability~\cite{chandrasekhar1961}, current sheets undergo tearing-mode reconnection~\cite{furth1963}, and magnetic flux tubes develop helical kink modes~\cite{bateman1978}.

In most regimes, analytical solutions are unavailable; progress hence depends on numerical simulation. Because plasma dynamics are multiscale, fine-scale structures usually coexist with large-scale flows, and uniform resolution is prohibitive. Direct Numerical Simulation (DNS) grid-point count scales as $Re^{9/4}$ in three dimensions~\cite{pope2000}, rapidly becoming intractable for high Reynolds numbers. Classical Adaptive Mesh Refinement (AMR)~\cite{berger1984,berger1989} increases resolution only where strong gradients or discontinuities are detected. However, current AMR refinement criteria are local and reactive. Their efficiency is limited if they respond only after instabilities have reached detectable amplitude.

Quantum computing offers alternative paradigms based on superposition, interference, and entanglement~\cite{nielsen2010}. Quantum approaches to linear algebra (HHL algorithm~\cite{harrow2009}) and quantum many-body simulation~\cite{feynman1982} are well studied, but direct quantum solution of nonlinear PDEs remains challenging: quantum circuits evolve linearly and unitarily, whereas MHD systems are nonlinear. Rather than forcing the two mathematical frameworks to reconcile, we shift our focus toward AMR as a strategic alternative to direct integration of the linear quantum framework.

The Quantum Approximate Optimisation Algorithm (QAOA)~\cite{farhi2014} alternates between a ``cost'' unitary (encoding problem structure) and a ``mixer'' unitary (generating quantum fluctuations), with variational parameters optimised by a classical loop. Its fixed-depth circuit structure and noise resilience make it suitable for Noisy Intermediate-Scale Quantum (NISQ) devices~\cite{preskill2018}.

This project develops a quantum-assisted steering mechanism for AMR: the Quantum-Hierarchical Adaptive Steering framework (Q-HAS). The quantum component serves a supportive role rather than replacing the classical solver. The minimisation of a structured (quantum Ising) Hamiltonian defined on a coarse-grained flux representation of the plasma domain acts as an instability detector. 
A major section of the following report presents an encoding of the physical quantities of plasma simulation into qubit states and the Hamiltonian.


The objectives of the project are threefold:
\begin{enumerate}
    \item To construct a mapping between coarse-grained MHD flux variables and qubit representations.
    \item To define and analyse the Hamiltonian governing instability detection.
    \item To implement the framework in software and assess, both theoretically and computationally, how such a hybrid mechanism performs compared to a typical classical AMR.
\end{enumerate}

The complete source code is publicly available at \url{https://github.com/armandld/BA_Proj}.

\chapter{Problem Formulation: The Flux Graph}
\label{chap:flux-graph}

Q-HAS represents the plasma domain as a graph in which each edge contains information about the physical fluxes. This chapter defines how the fluxes are represented in the edges and how classical plasma quantities are represented as qubit states.

\section{The Dual Graph Representation}

Let $\Omega \subset \mathbb{R}^3$ be the plasma volume. We decompose $\Omega$ into a set of coarse blocks $\mathcal{B} = \{b_1, \dots, b_N\}$.
Instead of nodal values like most classical methods, we define the state of the system by the flux exchange between these blocks.

Let $G = (V, E)$ be the graph where:
\begin{itemize}
    \item $V$ represents the volume blocks (nodes).
    \item $E$ represents the interfaces (surfaces or fluxes) between adjacent blocks.
\end{itemize}

We map the physical flux $\Phi_{ij}$ through each edge $(i,j) \in E$ to a qubit state.

\section{Classical Flux Extraction}

Let $U_i(t)$ and $U_j(t)$ be the classical state vectors at nodes $i$ and $j$ at time $t$. The implementation extracts five normalised MHD fields (velocity, magnetic and current fields): $\tilde{v}_x = v_x / v_0$, $\tilde{v}_y = v_y / v_0$, $\tilde{B}_x = B_x / B_0$, $\tilde{B}_y = B_y / B_0$, and $\tilde{J}_z = J_z / J_0$, where $v_0$, $B_0$, and $J_0$ are characteristic reference scales. For coarse blocks, the maximum value of these fields within the block is used.

For each edge (in 2D in this case), the differences between neighbouring cells are calculated using a shock-diode logic that retains only the physically significant contributions resulting from compression and shear (right/left for horizontal edges, up/down for vertical edges):

\begin{itemize}
    \item \textbf{Shock component (compression only):} Only negative normal differences (compression) are retained: $\Phi_{\text{shock}} = \sqrt{\max(0, -\Delta v_n)^2 + \max(0, -\Delta B_n)^2}$. Expansion (positive $\Delta v_n$) is filtered out as physically "harmless".
    \item \textbf{Shear component (always dangerous):} Tangential differences are taken in absolute value regardless of sign: $\Phi_{\text{shear}} = \sqrt{|\Delta v_t|^2 + |\Delta B_t|^2}$. Sliding motion is unstable in either direction.
    \item \textbf{Weighted combination:} The total stress flux combines both contributions with tuneable weights:
    \begin{equation}
        \Phi_{ij,l}(t) = \sqrt{(w_{\text{shock}} \cdot \Phi_{\text{shock}})^2 + (w_{\text{shear}} \cdot \Phi_{\text{shear}})^2 + (w_J \cdot \Delta J_z)^2}
    \end{equation}
\end{itemize}
\noindent where $l$ is for vertical or horizontal in 2D.
We retain the maximum risk from either side: $\Phi_{ij} = \max(\Phi_{ij}^{\text{right}}, \Phi_{ij}^{\text{left}})$.

The time derivative $\dot{\Phi}_{ij}(t)$ is the flux difference between two consecutive VQA updates:
\begin{equation}
    \dot{\Phi}_{ij, l}(t) = \Phi_{ij, l}(t) - \Phi_{ij, l}(t - \Delta t_{\text{hybrid}})
\end{equation}

\section{Embedding: Amplitude and Phase Encoding}
\label{sec:dense-embedding}

Physical observables are mapped to the full Bloch sphere coordinates $(\theta, \phi)$. Every single-qubit state is a point on the Bloch sphere~\cite{nielsen2010}, parameterised by a polar angle $\theta \in [0, \pi]$ and an azimuthal angle $\phi \in [0, 2\pi)$.

The polar angle $\theta$ is derived from a multi-indicator classical score, not from the stress flux $\Phi$ directly. This ensures each qubit starts with exactly the same initial probability as the classical AMR detector, so that the QAOA refines this baseline instead of replacing it completely. The stress flux $\Phi$ is used for the phase encoding $\phi$ and for the Hamiltonian coefficients (Chapter~\ref{chap:hamiltonian}), but not for $\theta$.

\subsection{Classical Multi-Indicator Score}
\label{sec:classical-score}

The amplitude encoding uses the same multi-indicator detection criteria as production MHD codes (FLASH~\cite{fryxell2000}, Athena~\cite{stone2008}, PLUTO~\cite{mignone2007}). Four MHD-AMR indicators are combined:
\begin{enumerate}
    \item \textbf{Vorticity:} $|\omega_z| = |\partial_x v_y - \partial_y v_x|$ \quad (shear layers, vortex cores, KH rolls)
    \item \textbf{Velocity divergence:} $|\nabla \cdot \mathbf{v}| = |\partial_x v_x + \partial_y v_y|$ \quad (compression, shocks)
    \item \textbf{Current density:} $|J_z|$ \quad (current sheets, reconnection sites)
    \item \textbf{L\"ohner estimator on $|\mathbf{B}|$:} A scale-free second-derivative error estimator~\cite{lohner1987}:
    \begin{equation}
        E_{\text{L\"oh}}(f) = \frac{\sqrt{(f''_{xx})^2 + (f''_{yy})^2}}{|f'_x| + |f'_y| + \epsilon_{\text{L\"oh}} |f| + \epsilon}
    \end{equation}
    which peaks at discontinuities and is insensitive to smooth gradients.
\end{enumerate}
Each indicator is normalised to $[0, 1]$ by its domain-wide maximum. The combined score is the root-mean-square (RMS) of the four normalised indicators:
\begin{equation}
    s(x,y) = \sqrt{\frac{s_{\text{vort}}^2 + s_{\text{div}}^2 + s_{J_z}^2 + s_{\text{L\"oh}}^2}{4}}
    \label{eq:classical-score}
\end{equation}
We used RMS instead of $\max$ to avoid over-refinement when a single noisy indicator fires. Usually, a cell needs multiple indicators to agree before the combined score triggers refinement.

\subsection{Amplitude Encoding ($\theta$)}

The polar angle $\theta$ is derived from the classical score $s$ such that the initial probability of measuring $\ket{1}$ equals the score exactly:
\begin{equation}
    \theta_{ij} = 2 \arcsin\!\left(\sqrt{s_{ij}}\right) \quad \Longrightarrow \quad P(\ket{1}) = \sin^2\!\left(\frac{\theta_{ij}}{2}\right) = s_{ij}
    \label{eq:theta-encoding}
\end{equation}
At initialisation (before any QAOA layers), each qubit has exactly the same probability of being flagged as the classical AMR detector. The QAOA refines this starting point using entanglement (plaquette and gradient topology in the Hamiltonian).

\subsection{Phase Encoding ($\phi$)}

The azimuthal phase $\phi$ encodes the temporal growth rate of the stress flux, enabling the VQA to detect constructive interference between adjacent cells whose instabilities grow simultaneously:
\begin{equation}
    \phi_{ij}(t) = \frac{\pi}{2} \cdot \tanh\!\left( \beta \frac{\Phi_{ij}(t) - \Phi_{ij}(t - \Delta t_{\text{hybrid}})}{\langle|\Delta\Phi|\rangle} \right)
    \label{eq:phi-encoding}
\end{equation}
Here, $\beta$ is a trainable hyperparameter controlling the sensitivity to rapid changes, and $\langle|\Delta\Phi|\rangle = \text{mean}(|\Phi(t) - \Phi(t-\Delta t_{\text{hybrid}})|)$ is the global average flux change, providing a dimensionless normalisation. This maps $\phi$ to the range $[-\frac{\pi}{2}, +\frac{\pi}{2}]$: $\phi \approx +\frac{\pi}{2}$ for locally growing instabilities, $\phi \approx 0$ for background evolution, and $\phi \approx -\frac{\pi}{2}$ for damping.

The quantum state of each edge qubit $\ket{q_{ij}}$ is initialised as:
\begin{equation}
    \ket{q_{ij}} = \cos\!\left(\frac{\theta_{ij}}{2}\right)\ket{0} + e^{i\phi_{ij}}\sin\!\left(\frac{\theta_{ij}}{2}\right)\ket{1}
\label{eq_init}
\end{equation}

Stable regions ($s \approx 0$) map to $\ket{0}$ (north pole), unstable regions ($s \approx 1$) to $\ket{1}$ (south pole), and the azimuthal phase captures the dynamical tendency.

\subsection{Gate-Level Implementation}
\label{sec:gate-level}

In the Qiskit implementation, the state~(\ref{eq_init}) is prepared using the general single-qubit rotation gate $R(\theta, \varphi)$, defined by:
\begin{equation}
    R(\theta, \varphi) = \cos\!\left(\frac{\theta}{2}\right) I - i \sin\!\left(\frac{\theta}{2}\right) \bigl(\cos(\varphi)\, X + \sin(\varphi)\, Y\bigr)
\end{equation}
Applying $R(\theta, \varphi)$ to $\ket{0}$ yields:
\begin{equation}
    R(\theta, \varphi)\ket{0} = \cos\!\left(\frac{\theta}{2}\right)\ket{0} + e^{i(\varphi - \pi/2)} \sin\!\left(\frac{\theta}{2}\right)\ket{1}
\end{equation}
To match the target state~(\ref{eq_init}) with azimuthal angle $\phi$, a phase offset of $\pi/2$ is applied: each qubit is initialised with the gate call $R(\theta_{ij}, \phi_{ij} + \pi/2)$, which ensures that the resulting state has the desired phase $e^{i\phi_{ij}}$ on the $\ket{1}$ amplitude.


\chapter{The Hamiltonian of Instability}
\label{chap:hamiltonian}

The Hamiltonian ground state identifies plasma regions that are topologically stressed or approaching instability. At high-Reynolds numbers, plasma dynamics usually arise from a finite number of physically distinct anomaly types: shear flows, vortices, magnetic reconnection, and helical kink modes. Each anomaly has a characteristic spatial signature encoded as a term in the quantum Hamiltonian.

Classical AMR detects instabilities only after they produce large gradients. Q-HAS encodes the structure of the plasma state at time $t$ to predict which regions will require refinement.

\section{Architecture Overview: Decoupled $f \times g \times T_{rcf}$}
\label{sec:hamilt-architecture}

The structure of the Hamiltonian is a decoupled multiplicative architecture. Each coefficient is a multiplicative product of three independent factors:
\begin{equation}
    \text{Coeff} = \text{Weight} \times g(\text{topology}) \times f(\text{scale}) \times T_{rcf}(\text{signal})
\end{equation}
It separates:
\begin{itemize}
    \item \textbf{Magnitude} of the anomaly ($T_{rcf}$: threshold-relative contrast filter),
    \item \textbf{Global thermodynamic scaling} ($f$: normal--critical gate),
    \item \textbf{Local topology} ($g$: leaky sigmoid gates).
\end{itemize}

The $ZZ$ and $ZZZZ$ interaction weights are fixed at 2 and 1 respectively (not trainable). The single-qubit $Z$ bias uses an adaptive weight $\alpha_z = w_{z,\text{frac}} \times \text{median}(|C|, |K|)$, where $w_{z,\text{frac}}$ is a trainable parameter. This keeps the $Z$ term subordinate to the spatial correlations ($ZZ$, $ZZZZ$) while breaking the QAOA ground-state degeneracy (without $Z$, the ferromagnetic $ZZ$/$ZZZZ$ coupling makes all-$\ket{0}$ and all-$\ket{1}$ degenerate).

The trainable parameters are: the encoding parameters ($\beta$), the uncertainty width ($\sigma$), the per-term contrast sensitivities ($\beta_{\text{curl}}$, $\beta_{\text{xpoint}}$), and the $Z$-bias fraction ($w_{z,\text{frac}}$), giving 5 trainable hyperparameters in total. Several parameters are frozen from prior calibration: $\text{threshold}_{\text{amr}} = 0.1496$ (from classical training), $\gamma_{\text{hydro}} = 2.0$, $\gamma_{\text{mag}} = 0.5$, and $\kappa = 10.0$. The $ZZ$ gradient coupling uses a fixed sensitivity $\beta = 1.0$ (its effective activation is controlled by $\sigma$ instead). All normalisation uses absolute physical non-dimensionalisation (dividing by fixed physical constants such as $Re_{\text{crit}} = 1$, $Rm_{\text{crit}} = 1$), never relative domain normalisation ($x / \max(x)$), ensuring that the Hamiltonian coefficients are stable across different grid resolutions and simulation times.

\section{The Hamiltonian Formula}

The Hamiltonian $H$ encodes the energy landscape of the plasma's topological stress. The VQA minimises $H$ to identify the configuration of stable and unstable regions most consistent with MHD anomaly physics.

QAOA~\cite{farhi2014} is used because: (1) the flux-graph formulation maps instability detection onto combinatorial graph optimisation (QAOA's native problem class), and (2) the Hamiltonian structure resembles the Ising model~\cite{lucas2014} (quantum Ising Hamiltonians are optimal computationally for QAOA because there is no need to diagonalise the Hamiltonian). QAOA alternates between two unitary operators:
\begin{itemize}
    \item The \textbf{structural (cost) Hamiltonian} $H_{\text{struct}}$ imprints a phase on each basis state proportional to its energy, marking configurations that correspond to physical anomalies.
    \item The \textbf{mixer Hamiltonian} $H_{\text{mixer}}$ generates quantum fluctuations (basis state transitions). Interference amplifies low-energy states and suppresses high-energy states.
\end{itemize}
Over $p$ layers, QAOA converges toward the ground state of $H_{\text{struct}}$.


\begin{equation}
    H_{\text{struct}} =
    \underbrace{\sum_i h_i Z_i}_{\substack{\text{Activity Bias} \\ \text{(adaptive } \alpha_z \text{)}}} +
    \underbrace{\sum_{\langle i,j \rangle} C_{ij} Z_i Z_j}_{\substack{\text{Gradient Coupling} \\ \text{(weight 2)}}} +
    \underbrace{\sum_{p} K_p \prod_{l \in p} Z_l}_{\substack{\text{Circulation Plaquette} \\ \text{(weight 1)}}} +
    \underbrace{\sum_{p} K_{\text{xpoint},p} \prod_{l \in p} Z_l}_{\substack{\text{X-point Reconnection} \\ \text{(weight 1, optional)}}}
\end{equation}

\noindent The X-point reconnection term detects hyperbolic magnetic null points via the determinant of the magnetic Jacobian (see Section~\ref{sec:xpoint-detection} below). It is a 4-body plaquette operator enabled by the \texttt{AdvAnomaliesEnable} flag. A fifth term for kink modes (Dzyaloshinskii--Moriya $XY - YX$ interactions driven by current helicity $\mathbf{J} \cdot \mathbf{B}$) is defined in the theoretical formulation but disabled in the 2D implementation (because kink modes are inherently three-dimensional, requiring helical deformation along the toroidal axis). The term is retained for future extension to 3D geometries. Note that $XY - YX$ is not diagonalizable in the $Z$ basis optimal for the QAOA algorithm (it is thus less optimal computationally to implement and run).

\noindent In this $H_{\text{struct}}$, each term acts as a specific ``sensor'' for a class of plasma instabilities, detailed below.

\begin{equation}
    H_{\text{mixer}} = \prod_i X_i
\end{equation}



\section{Activity Bias Term (The Validity Sensor)}
This single-body term biases each qubit according to whether its classical instability score exceeds the AMR threshold.
\begin{equation}
    H_{\text{bias}} = \sum_{i} h_i Z_i
\end{equation}
\textbf{Physical Coefficient $h_i$:} The bias coefficient directly encodes the deviation of the classical multi-indicator score (Equation~\ref{eq:classical-score}) from the AMR recursion threshold:
\begin{equation}
    h_i = \alpha_z \cdot \bigl(s_i - \text{threshold}_{\text{amr}}\bigr), \qquad \alpha_z = w_{z,\text{frac}} \times \text{median}\bigl(|C|, |K|\bigr)
\end{equation}
where $w_{z,\text{frac}}$ is trainable. The adaptive weight $\alpha_z$ scales with the median magnitude of $ZZ$ and $ZZZZ$ coefficients, keeping $Z$ subordinate to spatial correlations. Without $Z$, ferromagnetic coupling produces a degenerate ground state (all-$\ket{0}$ $=$ all-$\ket{1}$); the bias breaks this degeneracy. When $s_i > \text{threshold}_{\text{amr}}$, $h_i > 0$ biases toward $\ket{1}$ (refinement); when $s_i < \text{threshold}_{\text{amr}}$, $h_i < 0$ biases toward $\ket{0}$. Since $\theta$ is initialised from the same score $s$ (Equation~\ref{eq:theta-encoding}), the Z-bias and initial qubit amplitude are aligned, and the QAOA interaction terms can only flip qubits near the threshold where uncertainty lies.

\section{Gradient Coupling Term (The Boundary Sensor)}
This 2-body interaction term detects spatial discontinuities between neighbouring edges. A strong gradient indicates a boundary (shear layer, shock front, vortex boundary) that requires refinement.
\begin{equation}
    H_{\text{grad}} = \sum_{\langle ij, kl \rangle} C_{ij} Z_{ij} Z_{kl}
\end{equation}
\textbf{Physical Coupling $C_{ij}$:} The coupling follows the $f \times g \times T_{rcf}$ architecture, modulated by an uncertainty-weighted Gaussian centred at the AMR decision boundary:
\begin{equation}
    C_{ij} = 2 \cdot g_{\text{strain}}(Q_{\text{OW}}) \cdot \left\|\begin{pmatrix} f_{\text{hydro}}(Re) \cdot T_{\text{rcf}}(\Delta v) \\ f_{\text{mag}}(Rm) \cdot T_{\text{rcf}}(\Delta B) \end{pmatrix}\right\|
    \cdot \exp\!\left(-\left(\frac{\bar{s}_{ij} - \text{threshold}_{\text{amr}}}{\sigma}\right)^{\!2}\right)
\end{equation}
where $\bar{s}_{ij}$ is the average classical score of the two cells connected by the edge, the weight 2 is fixed (not trainable), and the $ZZ$ gradient terms use a fixed sensitivity $\beta = 1.0$ in their $T_{rcf}$ factor (sensitivity is controlled by $\sigma$ instead):
The Gaussian concentrates ZZ coupling near the decision boundary ($\bar{s} \approx \text{threshold}_{\text{amr}}$), where classical scoring is genuinely uncertain. Far from the threshold, ZZ coupling is exponentially suppressed and the Z-bias alone drives decisions. The trainable $\sigma \in [0.01, 0.30]$ controls the width of this activation band.
\begin{itemize}
    \item $g_{\text{strain}}(Q_{\text{OW}})$ is a leaky sigmoid topological gate based on the Okubo--Weiss Q-criterion~\cite{okubo1970, weiss1991} ($Q_{\text{OW}} > 0$ identifies strain-dominated regions):
    \begin{equation}
        g_{\text{strain}}(Q) = \frac{1}{1 + e^{-\kappa \, Q / Q_{\text{crit}}}}
    \end{equation}
    with $\kappa$ (trainable) controlling the steepness of the transition.
    \item $f(x)$ is a Normal-Critical gate with logarithmic growth above the turbulence threshold:
    \begin{equation}
        f(x, x_{\text{crit}}, \gamma) = \begin{cases} x / x_{\text{crit}} & x \leq x_{\text{crit}} \\ 1 + \gamma \ln(x / x_{\text{crit}}) & x > x_{\text{crit}} \end{cases}
    \end{equation}
    Applied separately for hydrodynamic ($f_{\text{hydro}}(Re)$) and magnetic ($f_{\text{mag}}(Rm)$) channels with $Re_{\text{crit}} = Rm_{\text{crit}} = 1$. The logarithmic form bounds growth ($Re = 3000$, $\gamma = 2$ gives $f \approx 17$ rather than $\infty$), preventing any single edge from dominating the Hamiltonian.
    \item $T_{rcf}(\text{val}, \text{val}_{\text{crit}}, \beta)$ is a threshold-relative contrast filter:
    \begin{equation}
        T_{rcf}(\text{val}, \text{val}_{\text{crit}}, \beta) = \beta \cdot \max\!\left(0,\; \frac{\text{val}}{\text{val}_{\text{crit}}} - 1\right)
    \end{equation}
    Unlike Michelson contrast (which kills the signal when the domain is uniformly active), $T_{rcf}$ compares to a fixed physical threshold, so the signal survives in pre-anomaly regimes. For $ZZ$, $\beta = 1.0$ is fixed (activation controlled by $\sigma$); for $ZZZZ$, $\beta_{\text{curl}}$ and $\beta_{\text{xpoint}}$ are trainable.
\end{itemize}
Depending on the characteristics of our physics flux, ($f \times g \times T_{rcf} \times \text{Gaussian}$) gives:
\begin{itemize}
    \item Below a fixed physical threshold ($\text{val} < \text{val}_{\text{crit}}$), $T_{rcf} = 0$ and $C_{ij} = 0$: no constraint is imposed.
    \item In regions where vortices dominate ($Q_{\text{OW}} < 0$), the strain gate suppresses $C_{ij}$ to $0$, preventing false gradient triggers inside vortex cores while maintaining differentiability.
    \item At strong strain boundaries ($Q_{\text{OW}} > 0$, high $Re$ or $Rm$, signal above threshold), all three factors are simultaneously large, producing a strong ferromagnetic coupling that enforces neighbour agreement on refinement decisions.
    \item Far from the decision boundary ($|\bar{s} - \text{threshold}_{\text{amr}}| \gg \sigma$), the Gaussian factor suppresses $C_{ij}$ exponentially, so the classical score alone drives decisions for clear-cut cells (it also helps reducing the depth of the quantum circuits). Near the boundary ($\bar{s} \approx \text{threshold}_{\text{amr}}$), the Gaussian factor is $\approx 1$ and ZZ correlations are fully active.
\end{itemize}

\section{Circulation Plaquette Term (The Vortex and Current Sensor)}
\label{sec:plaquettet}
Vorticity and current density are detected via 4-body interactions on closed loops (plaquettes), implementing a discrete Stokes theorem: the circulation $\oint \mathbf{F} \cdot d\mathbf{l}$ measures the enclosed curl.
\begin{equation}
    H_{\text{topology}} = \sum_{p \in \text{Plaquettes}} K_p \prod_{(ij) \in p} Z_{ij}
\end{equation}
\textbf{Physical Circulation via Tangential Projection:} Horizontal edges carry $v_x$ and $B_x$; vertical edges carry $v_y$ and $B_y$. This gives two independent physical curls:
\begin{align}
    \omega_z(i,j) &= v_x(i,j) - v_x(i,j{+}1) + v_y(i{+}1,j) - v_y(i,j) \approx (\partial_x v_y - \partial_y v_x) \cdot \Delta A \\
    J_{z,\text{curl}}(i,j) &= B_x(i,j) - B_x(i,j{+}1) + B_y(i{+}1,j) - B_y(i,j) \approx (\partial_x B_y - \partial_y B_x) \cdot \Delta A
\end{align}
where $\omega_z$ is the discrete vorticity, $J_{z,\text{curl}}$ is the discrete current density (Amp\`ere's law) and $\Delta A$ the "infinitesimal" area where the curl is evaluated. The plaquette coefficient follows the same decoupled $f \times g$ architecture as the gradient coupling, with independent topological gates for the fluid and magnetic components:
\begin{equation}
    K_p = \left\|\begin{pmatrix} g_{\text{rot}}(Q_{\text{OW}}) \cdot f_{\text{hydro}}(Re) \cdot T_{\text{rcf}}(\omega_z, \omega_{z,\text{crit}}, \beta_{\text{curl}}) \\ g_{\text{mag}}(|J_z|) \cdot f_{\text{mag}}(Rm) \cdot T_{\text{rcf}}(J_{z,\text{curl}}, J_{z,\text{crit}}, \beta_{\text{curl}}) \end{pmatrix}\right\|
\end{equation}
where the critical values are $\omega_{z,\text{crit}} = Re_{\text{crit}} \cdot \nu / dx^2$ and $J_{z,\text{crit}} = Rm_{\text{crit}} \cdot \eta / dx^2$, and:
\begin{itemize}
    \item $g_{\text{rot}}(Q_{\text{OW}})$ is a leaky sigmoid gate activated in rotation-dominated regions ($Q_{\text{OW}} < 0$), complementary to $g_{\text{strain}}$:
    \begin{equation}
        g_{\text{rot}}(Q) = \frac{1}{1 + e^{+\kappa \, Q / Q_{\text{crit}}}}
    \end{equation}
    \item $g_{\text{mag}}(|J_z|)$ is a leaky sigmoid gate activated when the current density exceeds a critical threshold:
    \begin{equation}
        g_{\text{mag}}(|J_z|) = \frac{1}{1 + e^{-\kappa \, (|J_z| / J_{\text{crit}} - 1)}}
    \end{equation}
    \item The $f$-gates and reference scales $v_0 = \langle |\mathbf{v}| \rangle$, $B_0 = \langle |\mathbf{B}| \rangle$ are as defined in the gradient coupling term.
\end{itemize}
This decoupling ensures that the fluid vorticity component is gated by $Q_{\text{OW}} < 0$ (rotation), while the magnetic current-sheet component is gated independently by $|J_z|$. A vortex without a current sheet activates only the $g_{\text{rot}}$ branch, and a current sheet without vorticity activates only the $g_{\text{mag}}$ branch.
\begin{itemize}
    \item $\omega_z$ detects vortices: regions of concentrated fluid rotation.
    \item $J_{z,\text{curl}}$ detects current sheets and magnetic reconnection sites.
    \item In uniform or irrotational regions, $K_p = 0$: no plaquette penalty is applied.
\end{itemize}

\section{X-point Reconnection \& Chirality Terms (Advanced Anomalies)}
\label{sec:xpoint-detection}
Additional Hamiltonian terms detect topological features of the magnetic field beyond standard turbulence.


\textbf{X-point Reconnection Detection (Plaquette Term, Optional):}
To detect magnetic reconnection at X-points (hyperbolic null points of the magnetic field), the Hamiltonian uses a 4-body plaquette operator based on the determinant of the magnetic Jacobian:
\begin{equation}
    H_{\text{xpoint}} = \sum_{p \in \text{Plaquettes}} K_{\text{xpoint},p} \prod_{l \in p} Z_l
\end{equation}
The coefficient is:
\begin{equation}
    K_{\text{xpoint}} = f_{\text{mag}}(Rm) \cdot T_{rcf}\!\bigl(\max(0,\; -\det(\nabla \mathbf{B})),\; \text{val}_{\text{crit}},\; \beta_{\text{xpoint}}\bigr)
\end{equation}
where $\det(\nabla \mathbf{B}) = (\partial_x B_x)(\partial_y B_y) - (\partial_x B_y)(\partial_y B_x)$ is the determinant of the magnetic field Jacobian. At an X-point, field lines form a hyperbolic pattern and $\det(\nabla \mathbf{B}) < 0$. The $\max(0, -\det)$ factor retains only topologically hyperbolic regions. The critical value for the threshold-contrast is $\text{val}_{\text{crit}} = (Rm_{\text{crit}} \cdot \eta / (dx \cdot B_0))^2$, and $\beta_{\text{xpoint}}$ is a trainable sensitivity specific to this term. This term is self-limiting: it activates only where the magnetic topology is genuinely hyperbolic, without requiring a separate topological gate. It is enabled by the \texttt{AdvAnomaliesEnable} flag and trained on the Island Coalescence scenario. Note : This allows the geometry to do most of the work. Instead of relying on rigid heuristics, the Jacobian determinant acts as a fairly natural ‘filter’ that ignores inactive magnetic islands and only comes into play when it detects a genuine reconnection site. I find this to be a mathematically elegant way of ensuring that the quantum solver correctly detects the reconnection points.


\section{Summary of Anomaly Mapping}
Table~\ref{tab:mapping} summarises how classical MHD phenomena are mapped to quantum operators in Q-HAS.

\begin{table}[h!]
\centering
\begin{tabular}{|l|l|l|l|c|}
\hline
\textbf{Anomaly} & \textbf{Physical Origin} & \textbf{Term} & \textbf{Operator} & \textbf{Weight} \\ \hline
Activity Bias & $\alpha_z \cdot (s_i - \text{threshold}_{\text{amr}})$ & Z bias & $Z_i$ & adaptive $\alpha_z$ \\ \hline
Shear / Gradient & $g_{\text{strain}} \times f \times T_{rcf} \times \text{Gauss}(\sigma)$ & ZZ coupling & $Z_i Z_j$ & 2 (fixed) \\ \hline
Vortex / Current & $\omega_z$, $J_{z,\text{curl}}$ (Stokes) & ZZZZ plaquette & $\prod_{\text{loop}} Z_k$ & 1 (fixed) \\ \hline
X-point Reconnection & $f_{Rm} \times T_{rcf}(-\det(\nabla\mathbf{B}))$ & ZZZZ plaquette & $\prod_{\text{loop}} Z_k$ & 1 (optional) \\ \hline
Kink (3D only) & $\mathbf{J} \cdot \mathbf{B}$ & DM interaction & $X_i Y_j - Y_i X_j$ & disabled \\ \hline
\end{tabular}
\caption{The Q-HAS Dictionary: Mapping MHD Instabilities to Quantum Operators. The $ZZ$ and $ZZZZ$ weights are fixed by operator degree. The $Z$ bias uses an adaptive weight subordinate to the spatial correlations.}
\label{tab:mapping}
\end{table}

\section{The Variational Ansatz}
The complete quantum circuit consists of two distinct stages: a data-embedding step followed by the QAOA evolution.

\begin{enumerate}
    \item \textbf{Data Embedding (Warm Start):}
    Each qubit is initialised from the physical flux data (Eq.~\ref{eq_init}), not from a uniform superposition. This prepares a product state $\ket{\psi_0}$ approximating the classical flux distribution:
    \begin{equation}
        \ket{\psi_0} \equiv \ket{\psi_0(\vec{\theta},\vec{\phi})} = \bigotimes_{\langle i,j \rangle} \left( \cos\!\left(\frac{\theta_{ij}}{2}\right)\ket{0} + e^{i\phi_{ij}}\sin\!\left(\frac{\theta_{ij}}{2}\right)\ket{1} \right)
    \end{equation}
    This warm-start strategy constrains VQA exploration to the neighbourhood of the physical solution, reducing QAOA iterations and mitigating the barren plateau problem~\cite{mcclean2018}.

    Note: Since the time evolution is encoded only in the ‘warm start’, a trade-off is inevitable: the longer the quantum circuit, the closer I get to minimising the Hamiltonian (the structures make perfect sense in the context of plasma physics). However, the ‘warm start’ becomes less and less significant, so that the time scale influences the Q-HAS probability map less and less. When the quantum circuit begins to be minimised, the phase $\phi$ guides the minimisation of the Hamiltonian (see \ref{sec:minimization-dynamics}).

    \item \textbf{QAOA Evolution:}
    The circuit applies $p$ layers of alternating unitary evolutions governed by $H_{\text{struct}}$ and $H_{\text{mixer}}$. For depth $p$, the trial state is:
    \begin{equation}
        |\psi(\vec{\Omega},\vec{\Gamma})\rangle = \prod_{k=1}^{p} \left( U_{\text{mixer}}(\Omega_k)\, U_{\text{struct}}(\Gamma_k) \right) |\psi_0(\vec{\theta},\vec{\phi})\rangle
    \end{equation}

    where:
    \begin{itemize}
        \item \textbf{The Phase Separator:} $U_{\text{struct}}(\Gamma_k) = e^{-i\Gamma_k H_{\text{struct}}}$ embeds the MHD constraints (shear, topology, reconnection) into the quantum phase of each basis state.
        \item \textbf{The Mixer:} $U_{\text{mixer}}(\Omega_k) = e^{-i\Omega_k H_{\text{mixer}}} = \prod_i e^{-i\Omega_k X_i}$ generates quantum fluctuations (spin flips) and facilitates interference between basis states, guided by the initial phases $\phi$.
    \end{itemize}

The variational parameters $\vec{\Gamma}, \vec{\Omega}$ are optimised by a classical loop (using the COBYLA optimiser~\cite{powell1994}) to minimise the energy expectation value $E = \langle \psi | H_{\text{struct}} | \psi \rangle$. The resulting ground state identifies the configuration of stable and unstable qubits that best fits the encoded physics.

\end{enumerate}



\section{Typical quantum circuit}

\begin{figure}[h]
\centering

\mbox{
    \Qcircuit @C=0.8em @R=1.2em {
        % ,  Line 1: Top Qubit , 
        \lstick{\ket{0}} 
        & \gate{R(\theta_{1,1,v},\phi_{1,1,v})} 
        & \multigate{3}{e^{i\Gamma_1H_{struct}}} % Couvre cette ligne + 3 en dessous (Total 4)
        & \multigate{3}{e^{i\Omega_1H_{mixer}}} 
        & \qw
        & \cdots 
        & {}
        & \qw
        & \multigate{3}{e^{i\Gamma_pH_{struct}}} 
        & \multigate{3}{e^{i\Omega_pH_{mixer}}} 
        & \qw
        & \meter \ar @{=>} [4,0]
        & \qw
        & \qw
        & \qw 
        \\
        % ,  Line 2: Middle Upper , 
        \lstick{\ket{0}} 
        & \gate{R(\theta_{1,1,h},\phi_{1,1,h})}  
        & \ghost{e^{i\Gamma_1H_{struct}}}
        & \ghost{e^{i\Omega_1H_{mixer}}}  
        & \qw
        & \cdots 
        & {}
        & \qw
        & \ghost{e^{i\Gamma_1H_{struct}}} 
        & \ghost{e^{i\Omega_pH_{mixer}}} 
        & \qw 
        & \qw
        & \meter \ar @{=>} [3,0]
        & \qw
        & \qw
        \\
        % ,  Line 3: Middle Lower , 
        \lstick{\vdots\hspace{0.2 cm}}
        & \vdots \hspace{0.1 cm}
        & {}
        & {} 
        & {}
        & \vdots 
        & {}
        & {}
        & {} 
        & {}
        & {} 
        & {}
        & {}
        & \ddots
        & {}
        \\
        % ,  Line 4: Bottom Qubit , 
        \lstick{\ket{0}} 
        & \gate{R(\theta_{n,n,h},\phi_{n,n,h})} 
        & \ghost{e^{i\Gamma_1H_{struct}}} 
        & \ghost{e^{i\Omega_1H_{mixer}}} 
        & \qw
        & \cdots 
        & {}
        & \qw
        & \ghost{e^{i\Gamma_1H_{struct}}} 
        & \ghost{e^{i\Omega_pH_{mixer}}} 
        & \qw
        & \qw
        & \qw
        & \qw
        & \meter \ar @{=>} [1,0] % Descend de 1 ligne vers le Bus
        \\
        % ,  Line 5: Classical Bus , 
        \lstick{Bus}
        & \cw {/}
        & \cw 
        & \cw 
        & \cw
        & \cdots 
        & \dstick{\raisebox{-0.3em}{\hspace{-1em}$\underbrace{\hspace{27.5em}}_{p \text{ layers}}$}}
        & \cw
        & \cw 
        & \cw
        & \cw
        & \cw
        & \cw
        & \cw
        & \cw \\
    }
}
\vspace{0.7cm}
\caption{Minimization of $\bra{\psi(\vec{\Omega},\vec{\Gamma})}H_{struct}\ket{\psi(\vec{\Omega},\vec{\Gamma})}$ with the QAOA circuit with Warm Start}
\end{figure}


\chapter{The Algorithm Structure}
\label{chap:algorithm}

Current quantum hardware has a limited number of available qubits, making full-grid QAOA infeasible. Our Q-HAS uses a hierarchical (multi-level) scanning strategy to overcome this implementation obstacle.

\section{Level-by-Level BFS Scanning}

The number of qubits is fixed, determining the patch size (at any level of execution): for a $\zeta$-dimensional periodic grid, $\zeta N_{\text{lim}}^{\zeta}$ qubits cover $N_{\text{lim}}^{\zeta}$ coarse blocks (the factor $\zeta$ counts the directions of the coarse block). All patches at the same depth are processed before moving to the next level.

At each depth level, physics is downsampled to the $N_{\text{lim}}^\zeta$ grid and used by the VQA to return a probability map (averaging horizontal and vertical edge marginals). For each sub-block:
\begin{itemize}
    \item If $\text{prob} \geq \text{threshold}_{\text{eff}}$: the sub-block is subdivided and added to the next depth level.
    \item If $\text{prob} < \text{threshold}_{\text{eff}}$: the sub-block is marked as a coarse leaf (calm region, solved at coarse resolution).
\end{itemize}

Recursion terminates at maximum depth $D_{\text{max}}$ or when the patch size falls below $\text{min\_size} = 6$ pixels. The termination depth determines the local solver resolution. The threshold is constant over depth (attempted to increase it linearly with depth at first to prevent excessive fine-scale refinement but a constant value over depth showed better results).\\

Note: The analysis depth (number of VQA levels evaluated) can be limited independently of the resolution depth (the resolution hierarchy used by the solver). For example, one might decide that if an instability is detected at an analysis depth of 3 (the third VQA level), then the region is computed at maximum resolution.

\section{Temporal Memory: Time-To-Live (TTL)}

An anomaly detected at time $t$ may weaken at $t + \Delta t_{\text{hybrid}}$ due to downsampling, noise, or transient fluctuation. Without temporal memory, the patch would be immediately dropped (false negative).

Each refined patch carries a \textbf{TTL (Time-To-Live) counter} (default $\tau = 1$ hybrid step). When flagged for refinement, TTL is set to $\tau$. At each subsequent step:
\begin{itemize}
    \item If the VQA probability is beyond the threshold: TTL is \textit{reset} to $\tau$ (fresh detection).
    \item If the VQA probability drops below the threshold but $\text{TTL} > 0$: the patch \textit{survives} and TTL is decremented by 1.
    \item If $\text{TTL} = 0$: the patch is released to the coarse level.
\end{itemize}
This stabilises the refinement tree: real instabilities have time to develop while transient false positives decay. The TTL map (keyed by patch bounds) persists across hybrid steps. Note: both VQA-guided and classical AMR paths maintain independent TTL maps.

\section{Boundary Probing}

After VQA, the engine checks whether the probabilities of the edge qubits (boundaries of the regions) are ‘too’ high. If an anomaly affects the boundary, the corresponding sub-cell is examined in greater detail, even in the case of a marginal probability (≥ 50\% of the effective threshold). This avoids the problem of “topology lost during zooming”.

The boundary detector compares the average probability of each edge row/column (top, bottom, left, right) to the inner average plus a tolerance. Active boundaries force the adjacent subcell to move to the next level of refinement.

\begin{algorithm}[H]
\caption{Q-HAS: Time-Stepping Loop}
\begin{algorithmic}[1]
\State \textbf{Initialize:} Classical Plasma State $\mathbf{U}_0$ on deep grid $G_0$ (resolution of $G_0$ is $N^{\zeta}$, with $\zeta$ the dimension of the simulation [2D or 3D])
\State Qubits number = $\zeta N_{lim}^{\zeta}$ 
\State $D_{max}= \lceil\frac{\log(N)}{log(N_{lim})}\rceil$
\State \textbf{Hyperparameters:} Convergence tolerance of the Quantum Approximate Optimization Algorithm (QAOA) $\epsilon$, Hamiltonian terms sensibilities and state sensibilities.
\For{$t = 0$ to $T_{end}$}
    \State \textit{// Step 1: Classical Coarse Update}
    
    \If{$t_{hybrid}|t$}
        \State \textit{// Step 1: Define coarse grid} 
        \State $G_{coarse}= \text{Reshape}(G_0, N_{lim})$
        \State \textit{// Step 2: Quantum Risk Assessment}
        \State Extract Fluxes $\Phi = F_{extract}(\mathbf{U}_{t},G_{coarse})$.
        \State Map Fluxes to Graph Parameters: $\mathbf{\theta, \phi} \leftarrow \mathcal{M}(\Phi)$.
        \State \textbf{VQA Routine:}
        \State \quad Minimize $\langle\psi(\vec{\Omega}, \vec{\Gamma})| H_{struct}|\psi(\vec{\Omega}, \vec{\Gamma})\rangle$ with QAOA.
        \State \quad Measure Qubits $\to$ Collapse to bitstring $z^*$.
        
        \State \textit{// Step 3: Identify High-Risk Zones}
        \State Let $\mathcal{R}$ be the set of edges where $z^*$ indicates instability (e.g., $|1\rangle$).
        
        \State \textit{// Step 4: Recursive Refinement (Zoom)}
        \If{$\mathcal{R}$ is not empty AND Depth $< D_{max}$}
            \State $\Omega_{sub} \leftarrow \text{Zoom}(G_0, \mathcal{R})$ \Comment{Define the spatial window in the high-resolution grid}
            \State $G_{coarse} \leftarrow \text{Reshape}(\Omega_{sub}, N_{lim})$ \Comment{Resample this area in $N_{lim}^d$ points}
            \State $\text{Depth} \leftarrow \text{Depth} + 1$
            \State \textbf{Recurse} Step 2 with the new grid $G_{coarse}$
        \EndIf
        \State Apply Fine Mesh AMR on final regions $G =\text{AMR}(G_0, \mathcal{R}) $.
    \EndIf
    
    \State \textit{// Step 5: Correction \& Advance}
    
    \State $\mathbf{U}_{t} \leftarrow \text{ClassicalSolver}(\mathbf{U}_t, G,  \Delta t_{coarse})$
\EndFor
\end{algorithmic}
\end{algorithm}

\chapter{Theoretical Advantage}
\label{chap:theory}
\section{Description of the Algorithm}

Q-HAS, like other AMR, aim to reduce the computational cost of the simulation by identifying which regions require high-resolution computation. Detection proceeds in two stages: (1)~flux amplitudes $\Phi_{ij}$ naively identify regions with strong gradients (as in classical AMR); (2)~the Hamiltonian $H_{\text{struct}}$ validates whether the pattern of flagged qubits is consistent with known MHD anomaly structures, filtering incoherent noise.

Two aspects of the algorithm provide the quantum advantage:
\begin{itemize}
    \item \textbf{Quantum tunnelling:} The QAOA mixer explores configurations a classical search could miss.
    \item \textbf{Phase-based anticipation:} Adjacent regions with high growth rate ($\dot{\Phi}$) but low amplitude ($\Phi$) are undetected by classical gradient detectors. In Q-HAS, by contrast, the aligned phases $\phi$ create a low-energy interference pattern, detecting pre-instability resonance before amplitude becomes critical.
\end{itemize}

The hierarchical structure provides an additional efficiency gain: determining the precise depth of each anomaly avoids unnecessary refinement in the surrounding regions.

\section{Initialization of Structural Parameters (of the QAOA)}
\label{sec:initialization}
\subsection{Initialization of $\vec{\Gamma}$}

While the Null-Gradient Proposition proves we cannot \textit{optimise} $\Gamma$ initially, we must still \textit{set} it. If $\Gamma=0$, the Mixer force $\nabla_\Omega$ also vanishes (as the state remains an eigenstate of $X$). Therefore, $\Gamma$ must be initialised to define the ``target phase'' that $\Omega$ will act upon.

We propose a Linear-Ramp ``Phase-Matching'' Strategy for general depth $p$.

\subsubsection{Optimal Initialization Rule}
To avoid phase aliasing ($\gamma E > 2\pi$ making high- and low-energy states indistinguishable), the total structural angle $\Gamma_{total} = \sum_{k=1}^p \Gamma_k$ must satisfy:
\begin{equation}
    \Gamma_{total} \approx \frac{\pi}{E_{max}}
\end{equation}

\begin{equation}
    E_{\max} \approx \sum_i |h_i| + \sum_{\langle i,j \rangle} |C_{ij}| + \sum_p |K_p|
\end{equation}

For a circuit of depth $p$, we distribute this total angle using a Linear Ramp Schedule (Trotterized Adiabatic Evolution):
\begin{equation}
    \Gamma_k^{init} = \left( \frac{2k}{p(p+1)} \right) \cdot \Gamma_{total} \quad \text{for } k=1 \dots p
\end{equation}
The ramp mimics Trotterized adiabatic evolution: early layers apply a small tilt (linear regime), later layers enforce constraints (locking in the solution as $H(t)$ transitions from $H_{mix}$ to $H_{struct}$).

\subsection{Initialization of $\vec{\Omega}$}

\subsubsection{The Null-Gradient Proposition (Structural Parameters)}
\label{sec:null-gradient}

\textbf{Proposition:} At the initialization point where all mixer parameters are set to zero ($\vec{\Omega} = \vec{0}$), the gradient of the energy with respect to any structural parameter $\Gamma_k$ is strictly zero, regardless of the circuit depth $p$.

\textbf{Proof:}
Let the depth be $p \geq 1$. The ansatz is:
\begin{equation}
    \ket{\psi(\vec{\Omega}, \vec{\Gamma})} = \prod_{k=1}^p \left( e^{-i\Omega_k H_{mixer}} e^{-i\Gamma_k H_{struct}} \right) \ket{\psi_0}
\end{equation}
At initialization, we set $\vec{\Omega} = \vec{0}$, so $e^{-i(0) H_{mixer}} = \mathbb{I}$. The state simplifies to a product of commuting exponentials:
\begin{equation}
    \ket{\psi(\vec{0}, \vec{\Gamma})} = \left( \prod_{k=1}^p e^{-i\Gamma_k H_{struct}} \right) \ket{\psi_0} = e^{-i (\sum_k \Gamma_k) H_{struct}} \ket{\psi_0}
\end{equation}
The energy expectation value is $E\equiv E(\vec{\Omega},\vec{\Gamma}) = \bra{\psi} H_{struct} \ket{\psi}$. Substituting the simplified state:
\begin{equation}
    E = \bra{\psi_0} e^{i (\sum \Gamma) H_{struct}} H_{struct} e^{-i (\sum \Gamma) H_{struct}} \ket{\psi_0}
\end{equation}
Since any operator commutes with itself ($[H_{struct}, e^{-i \alpha H_{struct}}] = 0$), the unitary operators cancel out:
\begin{equation}
    E = \bra{\psi_0} H_{struct} \ket{\psi_0} = \text{Constant}
\end{equation}
Consequently, the gradient is null:
\begin{equation}
    \nabla_{\Gamma_k} E \Big|_{\vec{\Omega}=0} = \frac{\partial}{\partial \Gamma_k} (\text{Const}) = 0 \quad \forall k \in [1, p]
\end{equation}
\textbf{Physical Implication:} The algorithm cannot lower energy by adjusting constraint weights alone ($\Omega = 0$). It must activate the Mixer ($\Omega$) to break symmetry. Phase-driven prediction is the exclusive engine of algorithm startup.

\section{Minimization Dynamics: Mathematical Proof of Advantage}
\label{sec:minimization-dynamics}
\subsection{Dynamics when $\vec{\Omega}= 0$ (initialization value)}
\subsubsection{The Driving Force (General Mixer Gradient)}
Since $\nabla_\Gamma E = 0$, optimisation is entirely driven by the mixer parameters $\vec{\Omega}$. The following generalises the gradient for any layer $k \in [1, p]$ at $\vec{\Omega} = \vec{0}$.

\textbf{Proposition:} The gradient $\partial_{\Omega_k} E$ is proportional to the commutator $[H_{\text{struct}}, H_{\text{mixer}}]$ evaluated on the phase-shifted state up to layer $k$.

\textbf{Proof:}
The unitary evolution after layer $k$ is $U_{>k}$. At $\vec{\Omega}=\vec{0}$, this consists solely of structural gates:
\begin{equation}
    U_{>k} = \prod_{j=k+1}^p e^{-i\Gamma_j H_{struct}} = e^{-i (\sum_{j=k+1}^p \Gamma_j) H_{struct}}
\end{equation}
The effective Hamiltonian at layer $k$ is $H_{eff} = U_{>k}^\dagger H_{struct} U_{>k}$. Since $[H, e^{i\alpha H}] = 0$, future layers cancel:
\begin{equation}
    H_{eff} = H_{struct}
\end{equation}
Thus, the gradient depends only on the state arriving at layer $k$, denoted $\ket{\psi_{k-1}}$:
\begin{equation}
    \partial_{\Omega_k} E \big|_{\vec{\Omega}=0} = -i \bra{\psi_{k}} [H_{struct}, H_{mixer}] \ket{\psi_{k}}
\label{eq:mixer-gradient}
\end{equation}
where $\ket{\psi_{k}} = \left( \prod_{j=1}^{k} e^{-i\Gamma_j H_{struct}} \right) \ket{\psi_0}$.\\

\noindent\textbf{Physical Implication:} The structural operator $e^{-i\Gamma_k H_{struct}}$ acts as a phase oracle. At $k=1$, $H_{struct}$ imprints topological defects into relative phases and $[H_{struct}, H_{mixer}]$ converts this into a non-zero gradient force. At $k>1$, the mixer exploits accumulated phase distributions from all preceding structural gates to resolve residual symmetries.


\subsubsection{Detailed Derivation by Hamiltonian Term}
Using $[Z, X] = 2iY$, the gradient for each term type is:

\paragraph{A. Gradient coupling ($ZZ$):} For $H_{shear} = C_{ij} Z_i Z_j$: $\nabla_{\Omega_m} E_{shear} \propto 2 C_{ij} \langle Y_i Z_j \rangle$.

\paragraph{B. Plaquette ($ZZZZ$):} For $H_{\mathcal{S}} = K \prod_{m \in \mathcal{S}} Z_m$ with qubit $k \in \mathcal{S}$:
\begin{equation}
    \nabla_{\Omega_m} E_{\mathcal{S}} \propto 2 K \left\langle \psi_m \left| Y_k \prod_{l \in \mathcal{S}, l \neq k} Z_l \right| \psi_m\right\rangle
\end{equation}

\subsubsection{From Gradient to Bit-Flip: The Physics of Detection}
\label{sec:coherence-filter}
How does a non-zero energy gradient transform a stable qubit ($\ket{0}$) into a flagged anomaly ($\ket{1}$)?

\paragraph{1. From Gradient to Bit-Flip}
The gradient update $\Omega_k^{(t+1)} = \Omega_k^{(t)} - \eta \cdot \partial_{\Omega_k} E$ is guaranteed non-trivial at initialisation (Equation~\ref{eq:mixer-gradient}). A negative gradient (instability detected) makes $\Omega_k$ positive, activating $U_{mixer}(\Omega) = e^{-i\Omega X}$: a Bloch sphere rotation from $\ket{0}$ (north pole) toward $\ket{1}$ (south pole), with $P(1) = \sin^2(\Omega)$.

\begin{tcolorbox}
\paragraph{3. The Coherence Filter (Why context matters)}
The magnitude of the rotation $\Omega$ is directly proportional to the gradient force $\vec{F}_q$ on qubit $q$:
\begin{equation}
    \vec{F}_q \propto \sin(\theta_q)\cdot\underbrace{\sin(\phi_q)}_{\text{Phase Boost}} \cdot\underbrace{\prod_{m \neq q} \cos(\theta_m)}_{\text{Topological Coherence}}
\label{eq:gradient-force}
\end{equation}

\textbf{Destructive interference (noise suppression):} If neighbours are incoherent ($\cos(\theta_m)$ averaging to 0), the product term vanishes, the gradient is zero, and the qubit remains at $\ket{0}$. The algorithm refuses to flag an isolated, topologically inconsistent anomaly.

\textbf{Constructive interference (structure detection):} If neighbours are laminar ($\cos\theta \approx 1$), the gradient is maximal ($\propto \sin(\phi_k)$), driving a large rotation $e^{-i\Omega X}$ that flips the qubit to $\ket{1}$ and flags the region for refinement.
\end{tcolorbox}
\subsection{Dynamics when $\vec{\Omega} \neq 0$ (Active Dynamics: The Competition of Torques)}

Once $\Omega_k > 0$, the null-gradient proposition for $\Gamma$ collapses. The system enters the active regime, governed by two competing unitary forces acting as orthogonal torques on the Bloch sphere.

\subsubsection{The Two Competing Forces}
The effective Hamiltonian at layer $k$ is $H_{eff}^{(k)} = U_{>k}^\dagger H_{struct} U_{>k}$, and the gradient with respect to parameter $\Lambda_k$ is $F_{\Lambda_k} = -i \bra{\psi_{k}} [G_k, H_{eff}^{(k)}] \ket{\psi_{k}}$. Two orthogonal torques govern the active dynamics:

\textbf{Interference force} ($\nabla_\Gamma$): acts on the Z-axis (phase rotation). With $\Omega \neq 0$, the effective Hamiltonian acquires off-diagonal terms, yielding $F_{\Gamma_k} = -i \bra{\psi_{k}} [H_{struct}, H_{eff}^{(k)}] \ket{\psi_{k}}$. The structural unitary imprints energy-dependent phases $|z\rangle \to e^{-i\Gamma E(z)}|z\rangle$; when the mixer couples basis states, high-energy configurations interfere destructively and low-energy configurations accumulate probability.

\textbf{Tunnelling force} ($\nabla_\Omega$): acts on the X-axis (bit-flip rotation). $F_{\Omega_k} = -i \bra{\psi_{k}} [H_{mixer}, H_{eff}^{(k)}] \ket{\psi_{k}}$ measures the torque exerted by problem constraints on the transverse axis. The mixer spreads the quantum state coherently across neighbouring configurations, redistributing amplitude beyond energy barriers (quantum tunnelling) and allowing escape from local minima.


\section{Algorithm Complexity Analysis}
\label{sec:complexity}

The computational cost per time step $\Delta t$ is compared across DNS, classical AMR, and Q-HAS.

\subsubsection{Direct Numerical Simulation (DNS)}
DNS solves on a uniform grid of resolution $N_{fine}^d$. The cost per step is $T_{\text{DNS}} = \mathcal{O}(N_{fine}^d)$; since $N_{fine} \sim Re^{3/4}$, this scales as $Re^{9/4}$ in 3D, making DNS intractable for fusion regimes.

\subsubsection{Classical Adaptive Mesh Refinement (AMR)}
Classical AMR refines only where a local gradient exceeds a threshold:
\begin{equation}
    T_{\text{AMR}} = \underbrace{\mathcal{O}(N_{fine}^d)}_{\text{Gradient Scan}} + \underbrace{\mathcal{O}(S_{class} \cdot N_{fine}^d)}_{\text{Physics Integration}} + \mathcal{O}(\text{Tree Overhead})
\end{equation}
where $S_{class}$ is the fraction of active cells. The bottleneck is reactive over-refinement: gradient triggers fire only after instabilities have grown, forcing large safety buffers that inflate $S_{\text{class}}$.

\subsubsection{Q-HAS Complexity Analysis}
While classical AMR is \textit{bottom-up} (scanning fine details), Q-HAS is \textit{top-down} (scanning coarse topological features to determine if fine details exist).

\paragraph{A. Complexity of the QAOA Subroutine (The Oracle)}
The VQA executes on fixed-size patches ($N_{\text{lim}} \times N_{\text{lim}}$). Execution time depends on optimisation steps $K_{\text{opt}}$ (bounded by $K_{\text{max}}$), shots $N_{\text{shots}}$, and circuit depth $p$:

\begin{equation}
    T_{\text{VQA}} = \mathcal{O}\left( \min(K_{max}, K_{opt}(\epsilon)) \cdot N_{shots} \cdot p \cdot N_{lim}^2 \right)
\end{equation}

Since $K_{max}$ is a fixed hyperparameter independent of the problem size, and $N_{shots}$ is constant, the oracle complexity is effectively $\mathcal{O}(1)$ with respect to the fine grid resolution $N_{fine}$.

While the cost scales polynomially with precision ($\sim \epsilon^{-3}$), it remains constant ($\mathcal{O}(1)$) with respect to the global simulation resolution $N_{fine}$.

\paragraph{B. Detection Overhead \& The Sparsity Advantage}
Even conservatively assuming a full gradient scan ($\mathcal{O}(N_{fine}^d)$) for flux extraction, the bottleneck in high-Reynolds MHD is numerical integration (Riemann solvers, magnetic induction in production codes), not error estimation. Performance gains are driven by the sparsity factor $S$: Q-HAS acts as a topological filter, minimising the expensive physics workload ($S_{\text{quant}} \ll S_{\text{class}}$). The total Q-HAS complexity simplifies (in worst case) to:


\begin{equation}
    T_{\text{Q-HAS}} \approx \underbrace{\mathcal{O}\left(\sum_{j=0}^{D_{max}} \left[ S_{quant} \cdot k^j \cdot T_{\text{VQA}} \right]\right)}_{\text{Quantum Tree Search}} + \underbrace{\mathcal{O}\left( S_{quant} \cdot N_{fine}^d \right)}_{\text{Sparse Physics Integration}}
\end{equation}

\begin{equation}
    T_{\text{Q-HAS}} \approx \underbrace{\mathcal{O}\left(\left[ S_{quant} \cdot k^{D_{max}} \cdot T_{\text{VQA}} \right]\right)}_{\text{Quantum Tree Search}} + \underbrace{\mathcal{O}\left( S_{quant} \cdot N_{fine}^d \right)}_{\text{Sparse Physics Integration}}
\end{equation}

\begin{equation}
    T_{\text{Q-HAS}} \approx \underbrace{\mathcal{O}\left( S_{quant} \cdot N_{fine}^d \cdot T_{\text{VQA}} \right)}_{\text{Quantum Tree Search}} + \underbrace{\mathcal{O}\left( S_{quant} \cdot N_{fine}^d \right)}_{\text{Sparse Physics Integration}}
\end{equation}

\noindent where $k^{D_{max}} = N_{fine}^d$ because $k = N_{lim}^d$ and $D_{max}=\frac{ln(N_{fine})}{ln(N_{lim})}$.
\subsubsection{The Spatio-Temporal Complexity Trade-off}
Let $\nu_{scan}$ ($0 < \nu_{scan} \le 1$) be how often the quantum detection layer is activated relative to physics steps. The amortised cost per step is:
\begin{equation}
    T_{\text{Q-HAS}}^{\text{avg}} \approx \mathcal{O}\left(S_{quant} \cdot N_{fine}^d \cdot (T_{\text{physics}} + \nu_{scan} \cdot T_{\text{VQA}})\right)
\end{equation}

The condition for quantum advantage over classical AMR is:
\begin{equation}
    \frac{S_{quant}}{S_{class}} < \frac{T_{\text{physics}}}{T_{\text{physics}} + \nu_{scan} \cdot T_{\text{VQA}}}
\end{equation}

This has 2 advantages: \textbf{spatial} ($S_{quant} \ll S_{class}$, topological filtering reduces the active volume) and \textbf{temporal} ($\nu_{scan} \ll 1$, because phase encoding detects the onset of instability, allowing refinement to persist for $\tau = 1/\nu_{scan}$ steps without re-querying the oracle).

\subsubsection{Future Scalability: The Boundary Sampling Regime}
With larger QPU devices ($N_{lim}$ large), Stokes’ theorem (already used by the plaquette operators in Section \ref{sec:plaquettet}) guarantees that rotational anomalies leave traces at the boundaries of the coarse blocks, thereby reducing the detection complexity from $\mathcal{O}(N_{fine}^d)$ to surface scaling $\mathcal{O}(N_{fine}^{d-1})$.


\chapter{Training Hyperparameters}
\label{chap:training}

Q-HAS relies on hyperparameters bridging classical plasma states and quantum operators. These parameters depend on the plasma regime ($Re$, $Rm$). Optimisation uses a \textit{Simulation-in-the-Loop} strategy with the Optuna Bayesian framework~\cite{akiba2019}.

\section{The Multi-Objective Cost Function}
The optimisation seeks to minimise a hybrid loss function $\mathcal{J}$ that penalises both physical error and computational resource usage:
\begin{equation}
    \mathcal{J}(\Theta) = \frac{\mathcal{E}_{\text{phys}}(\Theta) + \lambda \cdot \mathcal{S}(\Theta)}{1 + \lambda}
    \label{eq:combined-score}
\end{equation}
where $\Theta$ represents the set of trainable hyperparameters:\\ $\{\beta, \text{threshold}_{\text{amr}}, \gamma_{\text{hydro}}, \gamma_{\text{mag}}, \sigma, \beta_{\text{curl}}, \beta_{\text{xpoint}}, \kappa, w_{z,\text{frac}}\}$.

\begin{itemize}
    \item \textbf{Physical Fidelity ($\mathcal{E}_{\text{phys}}$):} We quantify accuracy using a weighted Relative $L_2$ Norm across five critical MHD fields: fluid velocity ($v_x, v_y$), magnetic field ($B_x, B_y$), and current density ($J_z$). We apply a weight map that combines both current density and vorticity:
    \begin{equation}
        w(x,y) = 1 + 0.5 \cdot \frac{|J_z(x,y)| / \langle|J_z|\rangle + |\omega_z(x,y)| / \langle|\omega_z|\rangle}{2}
    \end{equation}
    This ensures that missing an instability (false negative) in high-current or high-vorticity zones is penalised more heavily than over-refining a quiet region. Using both indicators ensures correct penalisation even in scenarios where only one is active (e.g., Lamb--Oseen vortex has high $|\omega_z|$ but low $|J_z|$).

    \item \textbf{Computational Sparsity ($\mathcal{S}$):} This represents the average fraction of the grid refined by the algorithm across all hybrid steps: $\mathcal{S} = \text{avg\_pixels\_refined} / N^2$.

    \item \textbf{The Trade-off Parameter ($\lambda$):} Set to $\lambda = 0.4$ in the Phase~1 training runs used in this report. This value imposes a meaningful sparsity constraint while keeping physical accuracy as the primary optimisation target.
\end{itemize}

\section{Three-Phase Progressive Training}
\label{sec:multi-scenario-training}
The optimisation was designed as a progressive three-phase schedule. In the work reported here, only Phase~1 was executed with split sensitivities (5 free quantum parameters) (not enough time for the analysis of others); the threshold was frozen from classical training. A DNS reference trajectory is precomputed once per scenario and reused across all Optuna trials, halving per-trial compute cost by skipping the witness solver.

\begin{enumerate}
    \item \textbf{Phase 1 (executed). Gate calibration with split sensitivities.} Trains 5 quantum parameters ($\beta$, $\sigma$, $\beta_{\text{curl}}$, $\beta_{\text{xpoint}}$, $w_{z,\text{frac}}$) on four scenarios simultaneously, each targeting a specific anomaly type:
    \begin{itemize}
        \item \textbf{Kelvin--Helmholtz:} Pure shear layers without current sheets, which calibrates $g_{\text{strain}}$ and the gradient coupling.
        \item \textbf{Harris Tearing:} Current sheet with reconnection sites, which calibrates $g_{\text{mag}}$ and selective detection.
        \item \textbf{Orszag--Tang:} Mixed turbulence with all anomaly types, providing a composite benchmark.
        \item \textbf{MHD Rotor:} Rotating field configuration generating vorticity from $t = 0$.
    \end{itemize}
    The frozen parameters ($\text{threshold}_{\text{amr}} = 0.1496$ from classical training, $\gamma_{\text{hydro}} = 2.0$, $\gamma_{\text{mag}} = 0.5$, $\kappa = 10.0$) were calibrated in earlier training runs or set by physical reasoning. The composite loss is the average over the four scenarios:
    \begin{equation}
        \mathcal{J}_{\text{composite}} = \frac{1}{4}\left(\mathcal{J}_{\text{KH}} + \mathcal{J}_{\text{Tearing}} + \mathcal{J}_{\text{OT}} + \mathcal{J}_{\text{Rotor}}\right)
    \end{equation}

    \item \textbf{Phase 2 (defined, not executed). Multi-scenario generalisation.} Would train all parameters on an expanded scenario set including Lamb--Oseen vortex and island coalescence.

    \item \textbf{Phase 3 (defined, not executed). Validation and fine-tuning on Orszag--Tang.} Would validate generalisation to the combined Orszag--Tang flow at $Re = Rm = 800$.
\end{enumerate}

\subsection{Hyperparameter Reference}

\begin{table}[H]
\centering
\caption{Trainable hyperparameters and their search ranges. Five quantum parameters are trained; four are frozen from prior calibration.}
\label{tab:hyperparams}
\begin{tabular}{llccc}
\toprule
\textbf{Parameter} & \textbf{Description} & \textbf{Range} & \textbf{Best} & \textbf{fANOVA} \\
\midrule
$\beta_{\text{curl}}$    & $T_{rcf}$ sensitivity for plaquette $ZZZZ$ terms & $[0.1, 10.0]$  & 4.27 & 44\% \\
$\beta_{\text{xpoint}}$  & $T_{rcf}$ sensitivity for X-point $ZZZZ$ terms   & $[0.1, 10.0]$  & 2.39 & 21\% \\
$\sigma$                 & Uncertainty width for ZZ coupling Gaussian        & $[0.01, 0.30]$ & 0.023 & 18\% \\
$\beta$                  & Phase encoding gain ($\phi$ sensitivity)           & $[0.5, 10.0]$  & 9.94 & 11\% \\
$w_{z,\text{frac}}$      & Z-bias fraction ($\alpha_z$ scaling)               & $[1.0, 200.0]$ & 10.40 & 6\% \\
\midrule
\multicolumn{5}{l}{\textit{Frozen parameters}} \\
$\text{threshold}_{\text{amr}}$ & AMR recursion threshold (from classical training) & -- & 0.1496 & -- \\
$\gamma_{\text{hydro}}$  & $f$-gate log-growth for $Re$                      & -- & 2.0 & -- \\
$\gamma_{\text{mag}}$    & $f$-gate log-growth for $Rm$                      & -- & 0.5 & -- \\
$\kappa$                 & Leaky sigmoid steepness                            & -- & 10.0 & -- \\
\bottomrule
\end{tabular}
\end{table}

\textbf{Interpretation of fANOVA results ("percentage of importance of each parameter to minimize the loss function" :} $\beta_{\text{curl}}$ (44\%) dominates, which means that the main driver of QAOA lies in the detection of magnetic reconnection via the curl-based $ZZZZ$ terms. $w_{z,\text{frac}}$ at only 6\% confirms that $ZZ$ spatial correlations barely influence the training score at $2\times 2$ VQA resolution. The trained $\sigma = 0.023$ is very tight, concentrating ZZ coupling on a narrow band around the threshold.

% ==========================================
% CHAPTER: COMPUTATIONAL IMPLEMENTATION
% ==========================================
\chapter{Computational Implementation}
\label{chap:implementation}

\section{The Hybrid Software Stack}
The framework is organised into four layers, all implemented in Python 3.11:
\begin{itemize}
    \item \textbf{Classical MHD Solver} (\texttt{src/Simulation/}): Built on \texttt{NumPy}~\cite{numpy2020} and \texttt{SciPy}~\cite{scipy2020}. The solver handles FD4 spatial discretisation, RK4 with CFL-adaptive time stepping, and FFT-based divergence-free projection on a $[0, 2\pi]^2$ periodic domain. The AMR engine (\texttt{refinement.py}) supports both quantum-guided BFS and classical-baseline scanning, with \texttt{step\_full()} for DNS and \texttt{step\_layered()} for hierarchical AMR.
    \item \textbf{Physics-to-Quantum Mapping} (\texttt{PhysToAngle.py}, \texttt{HamiltParams.py}): The \texttt{AngleMapper} computes the classical multi-indicator score (Equation~\ref{eq:classical-score}) for $\theta$ encoding and the stress flux $\Phi$ (with shock-diode logic) for $\phi$ encoding. The \texttt{PhysicalMapper} computes Hamiltonian coefficients using the $f \times g \times T_{rcf}$ architecture with absolute non-dimensionalisation.
    \item \textbf{Quantum Routine} (\texttt{src/VQA/}): A five-stage pipeline: (1)~circuit construction via Qiskit \texttt{QAOAAnsatz} with warm-start encoding, supporting periodic (depth~0) and bounded (depth~$>0$) Hamiltonians (the latter replaces boundary qubits with classical $\langle Z\rangle = \cos(\theta)$ expectation values); (2)~transpilation for the target backend; (3)~COBYLA-based minimisation with linear-ramp initialisation (Section~\ref{sec:initialization}), using \texttt{EstimatorV2} for energy and \texttt{SamplerV2} for measurement (\texttt{L-BFGS-B} and \texttt{Powell} are also supported); (4)~bitstring-to-probability post-processing; (5)~a \texttt{VQARuntime} singleton caching primitives and ansätze. Circuits run on the \texttt{Qiskit}~\cite{qiskit2024} \texttt{Aer} statevector simulator; IBM \texttt{FakeFez} and real hardware are supported via \texttt{--mode hardware} with dynamical decoupling and gate twirling for error mitigation.
    \item \textbf{Training Orchestration} (\texttt{TrainHyperParam\_v2.py}): \texttt{Optuna}~\cite{akiba2019} manages Bayesian optimisation with \texttt{MedianPruner} for early termination. Distributed training via PostgreSQL, JournalFileStorage, or SQLite.
    \item \textbf{Validation Suite} (\texttt{tests/}): \texttt{pytest}-based tests for solver convergence, AMR consistency, and VQA detection across eight study cases.
\end{itemize}

\subsection{The AMR Solver}

The AMR solver (\texttt{step\_layered}) implements a three-phase defect-correction architecture:
\begin{enumerate}
    \item \textbf{Phase 1 ,  Global coarse step:} The entire domain is downsampled by the coarse factor ($\text{target\_dim}^{D_{\max}}$), evolved by one RK4 step at coarse resolution, and the delta is upsampled back to full resolution (using bicubic periodic interpolation). This provides a cheap baseline correction for the whole domain.
    \item \textbf{Phase 2 ,  Per-patch finer resolution:} For each active patch, the local region is extracted with padding (8~pixels for the FD4$\times$RK4 stencil), downsampled to the patch's local resolution (determined by its depth), and evolved by a local RK4 step. Deeper patches (higher depth) have smaller \texttt{local\_factor}, yielding higher resolution.
    \item \textbf{Phase 3 ,  Global stabilisation:} The FFT-based divergence-free projection is applied to the full domain (identical to \texttt{step\_full}), ensuring incompressibility.
\end{enumerate}
\textbf{Convergence guarantee:} When all patches are at $D_{\max}$ (where $\texttt{local\_factor} = 1$), every patch's local RK4 step runs at full DNS resolution. The result is mathematically identical to \texttt{step\_full()}, up to floating-point precision ($< 10^{-10}$), as verified by the test suite.

\section{Numerical Methods: FD4 and RK4}
Numerical schemes are chosen for AMR compatibility:
\begin{itemize}
    \item \textbf{Spatial Discretisation:} 4th-order centred \textbf{Finite Differences (FD4)}. The local stencil allows independent solving on arbitrary rectangular sub-patches, essential for patch-based AMR (unlike global spectral methods).

    The Finite Volume Method (FVM) would be preferable: it derives from integral conservation laws, maintaining conservation exactly across resolution changes. FD requires interpolation at AMR boundaries, introducing non-conservative errors. FVM also acts as a natural low-pass filter via surface integrals ($\oint_{\partial\Omega} \mathbf{F} \cdot d\mathbf{S}$), whereas FD is sensitive to single-pixel noise. A full FVM implementation was beyond scope (see Chapter~\ref{chap:limitations}).

    
    \item \textbf{Temporal Integration:} 4th-order \textbf{Runge--Kutta (RK4)}. Centred FD schemes are non-dissipative, so RK4's stability region (function $R_4(z)= \sum_{k=0}^4 z^k / k!$, stable for $|R_4(z)|<1$) must cover the imaginary axis (approximately $[-2.82\,i,\; 2.82\,i]$).

    This is essential because FD4 applied to advection produces purely imaginary eigenvalues, requiring a time integrator with imaginary-axis stability coverage that RK2 lacks.

    \item \textbf{Incompressibility:} Two divergence-free constraints must hold at every instant:
    \begin{itemize}
        \item $\nabla \cdot \mathbf{v} = 0$ (mass conservation, constant density).
        \item $\nabla \cdot \mathbf{B} = 0$ (Gauss’s law, no magnetic monopoles).
    \end{itemize}
    RK4 and FD4 do not preserve these automatically; truncation errors create non-physical divergence. The \textbf{Projection Method} (Chorin~\cite{chorin1968}) enforces these constraints: after each step, an FFT-based spectral Poisson solver projects $\mathbf{v}$ and $\mathbf{B}$ onto divergence-free subspaces.

    For velocity: $\nabla^2 \varphi = \nabla \cdot \vec{v}^*$ (where $\vec{v}^*$ is the uncorrected field), then $\vec{v}_{\text{new}} = \vec{v}^* - \nabla \varphi$. The same applies to $\mathbf{B}$. FFT-based solving has $\mathcal{O}(N\log N)$ complexity.
\end{itemize}

\section{Software Validation and Test Suite}
\label{sec:test-suite}

A multi-tier test suite validates solver integrity and VQA anomaly detection sensitivity.

\subsection{Numerical Convergence and Solver Integrity}

Five convergence tests validate the AMR engine: exact convergence (\texttt{step\_layered} matches \texttt{step\_full} to $< 10^{-10}$ at max depth), normalisation of the patch cost metric, physics hierarchy (L2 error monotonically decreasing with depth), the \texttt{compute\_local\_factor} helper, and the corrected \texttt{max\_depth} subdivision formula.

\subsection{VQA Anomaly Study Cases}
\label{sec:vqa-study-cases}

Eight study cases validate each Hamiltonian term against synthetic plasma states: shear detection ($ZZ$), circulation ($ZZZZ$), X-point reconnection ($ZZZZ$), phase encoding ($\phi$), coefficient sign verification (all positive for correct minimisation), energy minimisation proof ($E_{\text{QAOA}} < E_{\text{all-}|0\rangle}$), cross-anomaly isolation (ensuring sensor independence), and combined anomalies (verifying simultaneous activation).

\section{Physical Scenarios}
\label{sec:physical-scenarios}

The solver supports eight MHD instability scenarios: Orszag--Tang vortex (mixed turbulence benchmark), Kelvin--Helmholtz (shear-driven rolling vortices), Harris tearing (current sheet with reconnection), magnetic twist (``silent'' instability invisible to classical AMR), noisy uniform field (noise immunity test), double tearing (frustrated current sheet interaction), Lamb--Oseen vortex (isolated vorticity), island coalescence (merging magnetic islands with X-points), and MHD rotor (rotating field generating vorticity from $t = 0$). The first four serve as the primary training scenarios; the remainder are used for isolated-anomaly validation.

\chapter{Computational Results}
\label{chap:results}


This chapter presents the performance tests designed to determine whether Q-HAS offers a measurable advantage over a conventional AMR benchmark model. The evaluation is deliberately conservative: both pipelines share the identical classical MHD solver, the same hierarchical BFS refinement engine, the same TTL mechanism, and the same domain, resolution and random seeds. The only difference is the decision variable: the classical pipeline thresholds the four-indicator score $s$ of Equation~\ref{eq:classical-score}, while the quantum pipeline thresholds the QAOA-refined probability map $P(|q\rangle = |1\rangle)=|\langle q|1\rangle|^2$. Each method uses its \textit{own} optimally-trained decision threshold, calibrated independently during training. Physical parameters are $Re = 400$, $Rm = 400$, $\eta = 0.001$, $c_s = 1.0$, $B_{z,\text{guide}} = 0.1$.

All figures in this chapter were generated by the \texttt{figures\_code/} scripts in the public repository and saved under \texttt{figures/}. The numerical values quoted in the text are taken directly from the figure log files.

\section{Training Results}
\label{sec:training-results}

\begin{table}[H]
\centering
\caption{Training outcome. Q-HAS achieves a marginal 0.66\% advantage in composite loss.}
\label{tab:training-outcome}
\begin{tabular}{lccc}
\toprule
\textbf{Method} & \textbf{Best Score} & \textbf{Best Trial} & \textbf{Trials Run} \\
\midrule
Q-HAS     & \textbf{0.2134} & \#4  & 170+ \\
Classical & 0.2148          & \#42 & 125 \\
\midrule
\textbf{Advantage} & \textbf{0.66\%} & -- & -- \\
\bottomrule
\end{tabular}
\end{table}

The quantum advantage is real but marginal. Optuna explored the parameter space widely ($w_{z,\text{frac}}$ from 7 to 168) without strong convergence, suggesting the classical threshold is already near-optimal at this scale ($N = 256$, $Re = Rm = 400$, $2\times 2$ VQA grid).

\subsection{Pareto Frontiers}
\label{sec:pareto}

Figure~\ref{fig:pareto-all} shows the Pareto fronts of the quantum and classical trials in the $(\mathcal{E}_{\text{phys}}, \mathcal{S})$ plane for $\lambda = 0.4$ (balanced value).

\begin{figure}[H]
    \centering
    \includegraphics[width=0.92\textwidth]{figures/fig0_pareto_all_combined.png}
    \caption[Pareto frontier: all scenarios combined]{Combined Pareto frontier across the four Phase~1 scenarios. The operating point at $\lambda = 0.40$ sits in the knee of the Pareto curve where further cost reduction yields diminishing physics fidelity.}
    \label{fig:pareto-all}
\end{figure}

% ===========================
\section{Per-Scenario Comparison}
\label{sec:per-scenario-results}

All metrics are computed against a ground-truth (GT) refinement map: a cell ``should be refined'' if its GT error (AMR vs DNS) exceeds the domain-wide mean.

\begin{itemize}
    \item \textbf{Captured fraction:} $\text{Captured} = \sum_{i \in \text{refined}} e_i \big/ \sum_{i \in \text{all}} e_i$, the fraction of total GT error inside refined cells.

    \item \textbf{Precision:} $|\text{refined} \cap \text{GT+}| / |\text{refined}|$, fraction of refined cells that needed refinement.

    \item \textbf{Recall:} $|\text{refined} \cap \text{GT+}| / |\text{GT+}|$, fraction of cells needing refinement that were refined.

    \item \textbf{Compute ratio:} fraction of domain refined ($\propto$ refined pixels$/N^2$ (not equal)).

    \item \textbf{Compactness (P/A):} perimeter-to-area ratio of the refined-patch union (lower = more coherent).

    \item \textbf{Component density:} connected components / total refined patches (lower = fewer fragments).

    \item \textbf{IoU stability:} overlap between refined-patch masks at successive time steps (Fig.~\ref{fig:temporal}).

    \item \textbf{Depth-resolved captured fraction:} captured fraction evaluated per BFS depth (Fig.~\ref{fig:depth-analysis}).
\end{itemize}

\subsection{Comprehensive Comparison}

\begin{figure}[H]
    \centering
    \includegraphics[width=\textwidth]{figures/fig4_comprehensive_comparison.png}
    \caption{Per-scenario captured fraction, precision, recall, and compute ratio for Q-HAS vs.\ classical AMR.}
    \label{fig:comprehensive}
\end{figure}

Table~\ref{tab:per-scenario} summarises the numerical results.

\begin{table}[H]
\centering
\caption{Per-scenario comparison: Q-HAS vs Classical AMR ($N = 256$, 2$\times$2 VQA, 8 qubits).}
\label{tab:per-scenario}
\begin{tabular}{lcccccccc}
\toprule
 & \multicolumn{2}{c}{\textbf{Captured}} & \multicolumn{2}{c}{\textbf{Precision}} & \multicolumn{2}{c}{\textbf{Recall}} & \multicolumn{2}{c}{\textbf{Compute}} \\
\textbf{Scenario} & Q-HAS & Class. & Q-HAS & Class. & Q-HAS & Class. & Q-HAS & Class. \\
\midrule
KH      & 0.951 & \textbf{0.968} & \textbf{0.722} & 0.683 & 0.975 & \textbf{1.000} & \textbf{0.730} & 0.836 \\
Tearing & \textbf{0.992} & 0.984 & 1.000 & 1.000 & 0.842 & \textbf{0.865} & 0.805 & \textbf{0.625} \\
Rotor   & \textbf{0.975} & 0.971 & 0.939 & \textbf{0.983} & 0.679 & \textbf{0.834} & 0.316 & \textbf{0.273} \\
OT      & \textbf{0.754} & 0.729 & 0.728 & \textbf{0.777} & 0.308 & \textbf{0.441} & 0.677 & \textbf{0.621} \\
\bottomrule
\end{tabular}
\end{table}

\textbf{Kelvin--Helmholtz (best efficiency case for Q-HAS):} Q-HAS uses 12.7\% less compute (0.730 vs 0.836) with only a marginal loss in captured fraction ($-1.8\%$) and better precision ($+5.7\%$). The KH instability has a simple, elongated structure; Q-HAS avoids over-refinement in the bulk.

\textbf{Harris Tearing:} Q-HAS achieves a higher captured fraction (0.992 vs 0.984) and identical precision (1.000), but at a much higher compute cost (0.805 vs 0.625, $+28.8\%$). The QAOA over-refines around the current sheet, capturing marginally more error at the expense of efficiency.

\textbf{MHD Rotor:} Q-HAS captures slightly more error (0.975 vs 0.971) but has lower precision ($-4.5\%$) and recall deficit ($0.679$ vs $0.834 = -18.6\%$), with a higher compute ratio (0.316 vs 0.273). The rotor's circular geometry with strong gradients at the interface produces QAOA corrections that tend to over-refine.

\textbf{Orszag--Tang (complex turbulence):} Q-HAS achieves a higher captured fraction (0.754 vs 0.729, $+3.4\%$) but with lower precision ($-6.3\%$), much lower recall ($0.308$ vs $0.441 = -30.2\%$), and higher compute (0.677 vs 0.621). The multi-scale turbulent structure produces many borderline cells where QAOA corrections are inconsistent.

\textbf{Overall assessment:} Q-HAS shows a genuine efficiency advantage only on Kelvin--Helmholtz (less compute, competitive capture). On Tearing and OT it captures slightly more error but at higher compute cost. On Rotor and OT it sacrifices precision and recall. There is no consistent global quantum advantage across scenarios.

\subsection{Statistical Validation}

\begin{figure}[H]
    \centering
    \includegraphics[width=\textwidth]{figures/fig6_statistical_validation.png}
    \caption{Bootstrap confidence intervals and permutation tests for Q-HAS advantage (10 seeds).}
    \label{fig:statistical}
\end{figure}

\begin{table}[H]
\centering
\caption{Statistical significance of Q-HAS advantage (bootstrap CI + permutation test, 10 seeds).}
\label{tab:statistical}
\begin{tabular}{lccccc}
\toprule
\textbf{Scenario} & \textbf{$\Delta$ captured} & \textbf{95\% CI} & \textbf{$p$-value} & \textbf{Cohen's $d$} & \textbf{Direction} \\
\midrule
KH      & $+0.0040$ & $[+0.0003, +0.0077]$ & 0.026 & $+0.65$ & Q-HAS better \\
Tearing & $+0.0032$ & $[+0.0025, +0.0039]$ & 0.001 & $+2.48$ & Q-HAS better \\
Rotor   & $-0.0104$ & $[-0.0235, -0.0004]$ & - & $-0.53$ & Classical better \\
OT      & $-0.0963$ & $[-0.1283, -0.0596]$ & - & $-1.60$ & Classical better \\
\bottomrule
\end{tabular}
\end{table}

Note: the script (of graph A of Fig.(\ref{fig:statistical})) uses a one-sided permutation test that detects only improvements. Cases where the conventional method is significantly better (MHD Rotor and Orszag--Tang) are labelled as ``n.s.''\ instead of being identified as statistically proven performance. (Due to the high computational cost, this figure retains the labels from the original script, which defaults to 'n.s.' for any result not favoring Q-HAS.)

% ===========================
\section{QAOA Mechanism Analysis}

\subsection{Detailed QAOA Corrections}

\begin{figure}[H]
    \centering
    \includegraphics[width=\textwidth]{figures/fig5_qaoa_detailed_analysis.png}
    \caption[Detailed QAOA analysis]{Detailed QAOA analysis.  Left: ground-truth error map.  Centre: QAOA correction $\delta = P_{\text{QAOA}}(|1\rangle) - s_{\text{classical}}$ per quadrant at depths~0 and~1.  Right: Hamiltonian term magnitudes ($H(Z)$, $C(ZZ)$, $K(ZZZZ)$) by BFS depth (log scale).}
    \label{fig:qaoa-detail}
\end{figure}

Corrections are \textbf{non-zero} at every BFS node, proving the QAOA circuit is functional.  The corrections are structured (they vary by scenario, depth, and cell position) and are not random noise.  The right column of Fig.~\ref{fig:qaoa-detail} reveals which Hamiltonian term dominates the cost function, and the answer is scenario-dependent:

\begin{itemize}
    \item \textbf{KH} (depth~1): $H(Z) \approx 1.5 \times 10^4$ and $K(ZZZZ) \approx 1.4 \times 10^4$ are comparable, while $C(ZZ) \approx 3.2 \times 10^3$ is an order of magnitude smaller.  The single-qubit bias and plaquette terms dominate; the gradient coupling contributes little.
    \item \textbf{Tearing} (depth~1): $C(ZZ) \approx 5.2 \times 10^3$ dominates, while $H(Z) \approx 31$ and $K(ZZZZ) \approx 163$ are two orders of magnitude smaller.  The gradient coupling is the primary signal, consistent with the strong spatial gradients across the current sheet.
    \item \textbf{Rotor} (depth~1): $H(Z) \approx 5.6 \times 10^5$ dominates, with $K(ZZZZ) \approx 1.4 \times 10^5$ and $C(ZZ) \approx 5.3 \times 10^4$.  At deeper levels this hierarchy reverses: by depth~3, $K > C > H$.
    \item \textbf{OT} (depth~1): $K(ZZZZ) \approx 6.6 \times 10^5$ dominates massively, followed by $C(ZZ) \approx 6.0 \times 10^4$ and $H(Z) \approx 3.5 \times 10^3$.  The topological plaquette term controls the OT cost landscape, consistent with OT having the richest topological structure (multiple interacting vortices and current sheets).
\end{itemize}

The Hamiltonian composition thus matches the physics of each scenario: gradient-dominated for Tearing (thin current sheet), plaquette-dominated for OT (interacting vortices), and bias-dominated for Rotor (strong circular interface).  Despite this physically coherent decomposition, the QAOA corrections (centre column) remain small ($|\delta| < 0.3$), meaning the VQA shifts probabilities by at most 30 percentage points from the classical score.  Whether these shifts improve or degrade accuracy depends on the scenario, as explored in Section~\ref{sec:decision-flip}.

\subsection{Hamiltonian Design Visualisation}

\begin{figure}[H]
    \centering
    \includegraphics[width=\textwidth]{figures/fig11_hamiltonian_design.png}
    \caption[Hamiltonian design visualisation]{Hamiltonian design visualisation (one row per scenario).  Columns: (1)~classical score $s$; (2)~uncertainty-weighted ZZ coupling ($\sigma = 0.023$ vs always-on); (3)~uncertainty weight field; (4)~Z-bias coefficient $h_i$.}
    \label{fig:hamiltonian-design}
\end{figure}

For KH, the ZZ coupling is effectively zero ($\sim 10^{-6}$): the Gaussian uncertainty weight\\ $\exp(-((\bar{s} - 0.15)/0.023)^2)$ kills all coupling because KH scores are bimodal (cells are either clearly above or clearly below the threshold, never near it).  Only for Rotor and OT does the ZZ coupling activate meaningfully (6.0\% and 6.1\% retention, respectively; Fig.~\ref{fig:sigma-ablation}), with OT exhibiting structured coupling along vortex filaments where scores cluster near the threshold.  This explains why $\sigma = 0.023$ limits quantum advantage: for scenarios with bimodal score distributions (KH, Tearing), the multi-body terms in the Hamiltonian are effectively null.

\subsection{Sigma Ablation Study}

\begin{figure}[H]
    \centering
    \includegraphics[width=\textwidth]{figures/fig13_sigma_ablation.png}
    \caption[Uncertainty weighting ($\sigma$) ablation]{Uncertainty weighting ($\sigma$) ablation.  (A)~Distribution of $|s - \text{threshold}|$ per scenario.  (B)~Fraction of edges with active ZZ coupling vs $\sigma$.  (C)~ZZ coupling magnitude retained relative to an always-on baseline.  At $\sigma = 0.023$, only Rotor and OT retain non-zero ZZ coupling.}
    \label{fig:sigma-ablation}
\end{figure}

At the trained $\sigma = 0.023$, only 3--15\% of edges have active ZZ coupling, and KH and Tearing retain 0\% of the ZZ magnitude (their score distributions are bimodal with cells far from the threshold).  Rotor retains 6.0\% and OT retains 6.1\%.  A broader $\sigma$ (0.10--0.15) would activate 87--99\% of edges, but Optuna chose $\sigma = 0.023$ because it minimised the composite training loss.  This implies that activating more ZZ edges does not help at the current scale: the corrections hurt as often as they help, so the optimiser suppresses them.

% ===========================
\section{Decision-Level Diagnostics}

\subsection{Decision Flip Analysis}
\label{sec:decision-flip}

\begin{figure}[H]
    \centering
    \includegraphics[width=\textwidth]{figures/fig15_decision_flip_analysis.png}
    \caption[Decision flip analysis]{Decision flip analysis ($N = 256$, depth = 5).  Columns: (1)~classical score vs QAOA probability per BFS cell; (2)~score distribution relative to threshold; (3)~flip fraction by BFS depth; (4)~QAOA correction magnitude vs distance-to-threshold.}
    \label{fig:flip-analysis}
\end{figure}

\begin{table}[H]
\centering
\caption{Decision flip analysis ($N = 256$, depth = 5). Classical is more accurate on aggregate.}
\label{tab:flips}
\begin{tabular}{lccccc}
\toprule
\textbf{Scenario} & \textbf{Cells} & \textbf{Flips} & \textbf{Flip Rate} & \textbf{QAOA Correct} & \textbf{Classical Correct} \\
\midrule
KH      & 904 & 125 & 13.8\% & 61.3\% & 65.6\% \\
Tearing & 584 & 109 & 18.7\% & 77.9\% & 81.2\% \\
Rotor   & 240 & 76  & 31.7\% & 67.9\% & 72.1\% \\
OT      & 408 & 145 & 35.5\% & 68.6\% & 70.3\% \\
\midrule
\textbf{Total} & 2136 & 455 & \textbf{21.3\%} & 43.3\% & 46.6\% \\
\bottomrule
\end{tabular}
\end{table}

The flip rate (21.3\%) shows that QAOA corrections change decisions at a meaningful fraction of cells.  However, classical is more accurate by 71 cells overall.  Median ZZ energy at flipped cells is $\sim 8 \times 10^5\times$ higher than at non-flipped cells (0.0078 vs $\sim 10^{-8}$), and flipped cells have $1.37\times$ higher mean GT error.  The multi-body coupling terms concentrate corrections on high-error regions near the decision boundary, but the corrections are not consistently in the correct direction.

\subsection{Decision Landscape}

\begin{figure}[H]
    \centering
    \includegraphics[width=\textwidth]{figures/fig16_decision_landscape.png}
    \caption[Decision landscape]{Decision landscape: classical score ($x$-axis) vs QAOA probability ($y$-axis) per BFS cell.  Green = should be refined (GT error above mean); brown = should not.  Dashed lines: decision threshold (0.150).  Off-diagonal cells are flips where the two methods disagree.}
    \label{fig:decision-landscape}
\end{figure}

\begin{table}[H]
\centering
\caption{Decision landscape: flip counts and net winner per scenario.}
\label{tab:landscape}
\begin{tabular}{lcccccc}
\toprule
\textbf{Scenario} & \textbf{Cells} & \textbf{Flips} & \textbf{Flip Rate} & \textbf{QAOA Correct} & \textbf{Class.\ Correct} & \textbf{Net Winner} \\
\midrule
KH      & 884 & 137 & 15.5\% & 61.4\% & 64.9\% & Classical by 31 \\
Tearing & 588 & 111 & 18.9\% & 77.7\% & 81.0\% & Classical by 19 \\
Rotor   & 252 & 80  & 31.7\% & 70.2\% & 75.0\% & Classical by 12 \\
OT      & 432 & 144 & 33.3\% & 69.9\% & 67.6\% & \textbf{Q-HAS by 10} \\
\bottomrule
\end{tabular}
\end{table}

OT is the \textit{only} scenario where Q-HAS wins on decision accuracy, consistent with OT having the most active ZZ/ZZZZ coupling and the richest topological structure.

% ===========================
\section{Topological Attribution}
\label{sec:topological-attribution}

\begin{figure}[H]
    \centering
    \includegraphics[width=\textwidth]{figures/fig17_topological_attribution.png}
    \caption[Topological correction attribution]{Topological correction attribution.  (A)~Decision accuracy binned by ZZ coupling strength quartile.  (B)~Flip accuracy by dominant Hamiltonian term (ZZ, ZZZZ, Z).  (C)~Q-HAS accuracy advantage in topology-rich vs topology-smooth sub-regions per scenario.}
    \label{fig:topological}
\end{figure}

\begin{table}[H]
\centering
\caption{Flip accuracy by dominant Hamiltonian term. ZZ-driven flips are 76.5\% accurate for Tearing; ZZZZ-driven flips are 66.3\% accurate for OT.}
\label{tab:flip-by-term}
\begin{tabular}{lcccccc}
\toprule
\textbf{Scenario} & \multicolumn{2}{c}{\textbf{ZZ (spatial)}} & \multicolumn{2}{c}{\textbf{ZZZZ (topology)}} & \multicolumn{2}{c}{\textbf{Z (bias)}} \\
 & Acc. & Flips & Acc. & Flips & Acc. & Flips \\
\midrule
KH      & 0.333 & 45 & 0.500 & 14 & 0.431 & 58 \\
Tearing & \textbf{0.765} & 68 & 0.000 & 8  & 0.000 & 41 \\
Rotor   & 0.429 & 14 & 0.491 & 53 & 0.333 & 15 \\
OT      & 0.378 & 45 & \textbf{0.663} & 92 & 0.200 & 5  \\
\bottomrule
\end{tabular}
\end{table}

\begin{table}[H]
\centering
\caption{Q-HAS advantage by topological regime. Topology-rich regions show positive advantage for Tearing ($+3.8$pp) and OT ($+5.3$pp).}
\label{tab:topo-regime}
\begin{tabular}{lcccc}
\toprule
\textbf{Scenario} & \textbf{Topology-rich $\Delta$} & \textbf{Smooth $\Delta$} & \textbf{Q-HAS rich} & \textbf{Classical rich} \\
\midrule
KH      & $-0.042$ & $-0.012$ & 0.322 & 0.363 \\
Tearing & $\mathbf{+0.038}$ & $-0.082$ & 0.729 & 0.692 \\
Rotor   & $-0.057$ & $-0.008$ & 0.631 & 0.689 \\
OT      & $\mathbf{+0.053}$ & $+0.024$ & 0.675 & 0.621 \\
\bottomrule
\end{tabular}
\end{table}

\textbf{Central finding:} In topology-rich regions, Q-HAS shows a positive advantage for Tearing ($+3.8$ percentage points) and OT ($+5.3$pp). These are the two scenarios with the strongest topological features (current-sheet reconnection for Tearing, vortex interactions for OT). In smooth regions, the advantage disappears or reverses.

\textit{Note} (Harris Tearing): The decision flip analysis (Fig.~\ref{fig:flip-analysis}) shows 109 flips on Tearing, of which 45 are QAOA-correct and 64 are classically-correct.  When decomposed by dominant Hamiltonian term (Fig.~\ref{fig:topological}), ZZ-driven flips are 76.5\% accurate (52/68), but Z-bias-driven flips are 0\% accurate (0/41) and ZZZZ-driven flips are 0\% accurate (0/8).  The gradient coupling ($ZZ$) is well-calibrated for this scenario, but the bias ($Z$) and plaquette ($ZZZZ$) terms produce uniformly incorrect corrections.

% ===========================
\section{Physical Fidelity and Structural Analysis}

\subsection{Physical Fidelity}

Figure~\ref{fig:hierarchical-overview} gives an overview of the four key metrics across scenarios: captured error fraction (panel A), total patch count (panel B), compute ratio (panel C), and L2 error versus DNS (panel D).  Q-HAS and classical AMR are closely matched on captured fraction ($>0.95$ for KH, Tearing, Rotor; $>0.72$ for OT).  Q-HAS uses more patches and compute on three of four scenarios (Tearing, Rotor, OT), with KH being the exception (0.730 vs 0.836).  Classical AMR achieves lower L2 error on all four scenarios (panel D).

\begin{figure}[H]
    \centering
    \includegraphics[width=0.75\textwidth]{figures/fig8_hierarchical_comparison.png}
    \caption[Hierarchical AMR overview]{Hierarchical AMR overview. (A)~Captured error fraction; (B)~patch count; (C)~compute ratio; (D)~L2 error vs DNS.}
    \label{fig:hierarchical-overview}
\end{figure}

Figure~\ref{fig:physical-fidelity} tracks the L2 error, kinetic energy, and enstrophy conservation over 15 AMR-guided evolution steps, probing whether the two methods diverge dynamically.

\begin{figure}[H]
    \centering
    \includegraphics[width=\textwidth]{figures/fig7_physical_fidelity.png}
    \caption[Physical fidelity over time]{L2 error vs DNS, kinetic energy, and enstrophy conservation over 15 AMR-guided steps.  All L2 errors $< 10^{-3}$; classical tracks slightly closer to DNS on every scenario.}
    \label{fig:physical-fidelity}
\end{figure}

\begin{table}[H]
\centering
\caption{Final L2 error vs DNS after 15 AMR steps ($N = 256$, depth = 5). Classical wins on all scenarios.}
\label{tab:fidelity}
\begin{tabular}{lccc}
\toprule
\textbf{Scenario} & \textbf{Q-HAS Final L2} & \textbf{Classical Final L2} & \textbf{Winner} \\
\midrule
KH      & $2 \times 10^{-6}$ & $\sim 0$           & Classical \\
Tearing & $1.0 \times 10^{-5}$ & $3 \times 10^{-6}$ & Classical \\
Rotor   & $9.6 \times 10^{-4}$ & $5.2 \times 10^{-4}$ & Classical \\
OT      & $1.6 \times 10^{-4}$ & $\sim 0$           & Classical \\
\bottomrule
\end{tabular}
\end{table}

Both methods achieve excellent fidelity (all L2 errors $< 0.001$). The differences are practically negligible for KH and Tearing (L2 $< 10^{-5}$). Energy and enstrophy conservation track closely for both methods.

\subsection{Spatial Coherence}

\begin{figure}[H]
    \centering
    \includegraphics[width=\textwidth]{figures/fig3_spatial_coherence.png}
    \caption[Spatial coherence analysis]{Spatial coherence analysis.  Left: Q-HAS (red) and classical (green) patch outlines.  Centre: compactness and component density per scenario.  Right: precision vs GT error percentile threshold.}
    \label{fig:spatial-coherence}
\end{figure}


\begin{table}[H]
\centering
\caption{Spatial coherence metrics (lower = more coherent). Classical AMR is consistently more spatially coherent.}
\label{tab:coherence}
\begin{tabular}{lcccc}
\toprule
 & \multicolumn{2}{c}{\textbf{Compactness (P/A)}} & \multicolumn{2}{c}{\textbf{Component Density}} \\
\textbf{Scenario} & Q-HAS & Classical & Q-HAS & Classical \\
\midrule
KH      & 0.050 & \textbf{0.040} & 0.06 & 0.06 \\
Tearing & 0.072 & \textbf{0.070} & 0.12 & 0.12 \\
Rotor   & 0.142 & \textbf{0.067} & 0.50 & \textbf{0.14} \\
OT      & 0.219 & \textbf{0.190} & 2.04 & \textbf{0.93} \\
\bottomrule
\end{tabular}
\end{table}

On KH and Tearing, compactness and component density are nearly identical between the two methods.  On Rotor and OT, however, Q-HAS produces noticeably more fragmented patch layouts (Rotor compactness $2.1\times$ worse; OT component density $2.2\times$ worse).  This is counterintuitive: the ZZ correlations were designed to enforce spatial consistency (ferromagnetic coupling), but on complex scenarios the perturbative corrections push some cells above threshold while pulling adjacent cells below, creating isolated patch fragments.  The ZZ coupling is too weak (suppressed by $\sigma = 0.023$) to impose coherence.

\subsection{Early Detection and Temporal Stability}

\begin{figure}[H]
    \centering
    \includegraphics[width=\textwidth]{figures/fig2_early_detection.png}
    \caption[Early detection and temporal stability]{Early detection and temporal stability.  Rows: recall, precision, and IoU stability over time.  Columns: one per scenario.}
    \label{fig:temporal}
\end{figure}

Q-HAS shows slightly higher precision than classical on KH ($+3.3$pp average) and OT ($+1.5$pp average), driven by the QAOA selectively removing some false-positive patches.  On the other hand, classical has higher recall on every scenario (it refines more aggressively) and exhibits more temporally stable patch layouts: the IoU curves for classical are nearly flat on KH and Tearing, while Q-HAS shows temporal fluctuations (IoU drops of up to 0.25 on Rotor) caused by stochastic QAOA corrections that change which patches cross the threshold from step to step.

\subsection{Depth Analysis}

\begin{figure}[H]
    \centering
    \includegraphics[width=\textwidth]{figures/fig12_depth_analysis.png}
    \caption[Depth-resolved AMR analysis]{Depth-resolved AMR analysis for a short time $\delta t$ (not a complete simulation).  Columns: (1)~patches per BFS depth; (2)~GT error captured per depth; (3)~decision agreement per depth; (4)~compute budget per depth.  GT error concentrates at depth~5; coarser levels serve as spatial routing.}
    \label{fig:depth-analysis}
\end{figure}

Through depths 0--4, Q-HAS and classical agree on 100\% of refinement decisions across all scenarios.  Disagreements appear only at depth~5 (the finest level), where agreement drops to 91.9\% on KH, 99.2\% on Tearing, 97.2\% on Rotor, and 92.6\% on OT.  This confirms that the QAOA corrections are perturbative: the hierarchical BFS structure is robust, and the quantum circuit modifies decisions only at the finest scale where the scoring landscape is most ambiguous.

% ===========================
\chapter{Discussion of Results}
\label{chap:discussion}

\section{Discussion of Theoretical Results}

The gradient force (equation~\eqref{eq:gradient-force}) linking the physical state of the plasma ($\theta$, $\phi$) to the direction of variational optimisation seems intuitively very consistent. Indeed, it is quite impressive to see the temporal prediction ($\sin(\phi_q)$, growth rate) align so well with the topological consistency ($\prod_{m \neq q} \cos(\theta_m)$, spatial context) in the expression of the force.\\

\noindent The proof that $\nabla_{\Gamma} E = 0$ at $\vec{\Omega} = 0$ (Section~\ref{sec:null-gradient}) shows that structural parameters alone cannot initiate optimisation; the mixer must activate first. This nicely justifies the linear-ramp strategy for $\vec{\Gamma}$.\\

\noindent Finally, the coherence filter (Section~\ref{sec:coherence-filter}) distinguishes genuine instabilities from noise. In incoherent regions (like noisy regions), the topological coherence factor averages to zero, suppressing the gradient and preventing false-positive refinement.

\section{Cross-Figure Synthesis}

\begin{enumerate}
    \item \textbf{The QAOA circuit is functional.} Fig.~\ref{fig:qaoa-detail} proves non-zero corrections ($\delta$) exist at every BFS node. The corrections are structured because they vary by scenario, depth, cell position (not random noise).

    \item \textbf{The corrections are perturbative, not transformative.} Captured fractions differ by at most 2.5 percentage points between methods (Fig.~\ref{fig:comprehensive}), and precision differences remain below 5pp.  However, recall differences are substantial on Rotor ($-15.5$pp) and OT ($-13.3$pp), indicating that the QAOA corrections selectively suppress refinement at the finest depth.

    \item \textbf{KH is the only scenario with a small efficiency advantage.} Q-HAS uses 12.7\% less compute with nearly the same capture quality, and this is statistically significant ($p = 0.026$). On Tearing, Q-HAS captures more error (0.992 vs 0.984) but at 28.8\% higher compute cost. On Rotor and OT, Q-HAS is clearly worse on recall and precision.

    \item \textbf{ZZ correlations fragment rather than consolidate patches on complex scenarios (Fig.~\ref{fig:spatial-coherence}).} On KH and Tearing the patch layouts are nearly identical between methods, but on Rotor and OT, Q-HAS produces measurably more fragmented patches (compactness $2.1\times$ worse on Rotor, component density $2.2\times$ worse on OT).  The coupling strength is in a regime where it perturbs individual cell decisions without being strong enough to enforce spatial coherence across the patch.

    \item \textbf{The root cause is scale?} At $2\times 2$ VQA resolution, only 4 cells are coupled, and classical scoring is already near-optimal at $Re = 400$.  Although the Hamiltonian composition is physically coherent (gradient-dominated for Tearing, plaquette-dominated for OT; Fig.~\ref{fig:qaoa-detail}), $\sigma = 0.023$ suppresses ZZ coupling on most cells (Fig.~\ref{fig:sigma-ablation}), and the resulting corrections are too small and inconsistent to yield a net global advantage.  Disagreements between methods appear only at the deepest BFS level (100\% agreement through depths 0--4, dropping to 92--99\% at depth~5; Fig.~\ref{fig:depth-analysis}).  OT is the only scenario where Q-HAS wins on decision accuracy (Table~\ref{tab:landscape}), consistent with it having the strongest plaquette ($ZZZZ$) signal. $\sigma$ could be interpreted either as a limit on the Q-HAS scale caused by constraints that restrict the identification of constraints (which would be lifted as $N_{lim}$ increases) or as indicating that Q-HAS decisions should follow the classical AMR's decisions as much as possible (disadvantage of Q-HAS).  

    \item \textbf{Topology-rich regions show measurable Q-HAS advantage (Fig.~\ref{fig:topological}).} ZZ-driven flips are 76.5\% accurate for Tearing; ZZZZ-driven flips are 66.3\% accurate for OT. In topology-rich sub-regions, Q-HAS has $+3.8$pp advantage for Tearing and $+5.3$pp for OT. This is the strongest evidence that the Hamiltonian's spatial coupling terms detect real physics, even though the overall advantage could be diluted because most of the domain is topologically smooth.
\end{enumerate}

\section{Identified Errors and Methodological Limitations}
\label{sec:identified-errors}

The following methodological limitations qualify the conclusions:

\begin{enumerate}

    \item \textbf{Term-specific failure modes (Fig.~\ref{fig:topological}).}
    On Harris Tearing, Q-HAS makes 64 incorrect flips vs 45 correct flips (Fig.~\ref{fig:flip-analysis}).  The per-term decomposition reveals that ZZ-driven flips are 76.5\% accurate (52/68), but Z-bias flips are 0\% accurate (0/41) and ZZZZ flips are 0\% accurate (0/8).  The gradient coupling term works as intended; the bias and plaquette terms do not.  On all four scenarios collectively, classical is correct on 71 more cells than Q-HAS (996 vs 925 out of 2136).

    \item \textbf{High classical bias weight (Fig.~\ref{fig:flip-analysis}).}
    The trained $w_{z,\text{frac}} = 10.40$ causes the Z-bias term to dominate the Hamiltonian on Rotor (Fig.~\ref{fig:qaoa-detail}).  At 78.7\% of cells, Q-HAS and classical produce the same decision; the QAOA circuit modifies decisions at only 21.3\% of cells.  This is consistent with the optimiser finding that, at $2\times 2$ resolution, the classical score is already a strong predictor.

    \item \textbf{Ground-truth definition.}
    The GT threshold (``should refine'' $=$ GT error $>$ mean GT error) is a methodological choice, not a physical constant.  All accuracy and flip-correctness metrics depend on this definition; a different threshold would change the numerical conclusions.
\end{enumerate}

\section{Potential of Q-HAS}

Three hypotheses, not yet tested, suggest where quantum advantage could emerge.  First, at higher VQA resolution ($4\times 4$ or $8\times 8$), ZZ correlations couple more cells; boundary correction potential scales as $\mathcal{O}(N)$ while total cells scale as $\mathcal{O}(N^2)$, so the relative influence of coupling grows.  Second, at higher Reynolds numbers ($Re/Rm \gg 800$), classical indicator scores become noisier (more borderline cells near the threshold), which could give ZZ and ZZZZ correlations more signal to exploit; retraining would be required.  Third, scenarios with richer topological structure could amplify the localised advantage measured in Fig.~\ref{fig:topological} ($+3.8$pp on Tearing, $+5.3$pp on OT in topology-rich sub-regions) into a global one.  These remain conjectures pending further experiments.

\chapter{Conclusions}
\label{chap:conclusions}


A consistent mapping from coarse-grained MHD observables to qubit states was constructed. The polar angle $\theta = 2 \arcsin(\sqrt{s})$ is derived from a classical multi-indicator score (RMS of vorticity, divergence, $|J_z|$, L\"ohner estimator), ensuring each qubit starts with the same probability as the classical AMR detector. The azimuthal phase $\phi$ encodes the stress flux growth rate. In principle, this implementation should allow QAOA to combine spatial indicators with temporal information on the growth rate; however, current experiments do not isolate or demonstrate earlier detection of instabilities than classical AMR.\\

A structured Hamiltonian was defined using the decoupled $f \times g \times T_{rcf}$ architecture, with individual terms corresponding to four physically distinct classes of MHD anomalies: activity bias (single-qubit $Z$ with adaptive weight $\alpha_z$), uncertainty-weighted spatial gradient coupling ($ZZ$ with weight 2, modulated by $\exp(-((s - \text{thr})/\sigma)^2)$ to focus on the decision boundary), physical circulation (plaquette $ZZZZ$ with weight 1), and X-point reconnection (plaquette $ZZZZ$ with weight 1, optional). Five parameters were trained ($\beta$, $\sigma$, $\beta_{\text{curl}}$, $\beta_{\text{xpoint}}$, $w_{z,\text{frac}}$); where four were already frozen and determined intuitively/arbitrarily ($\text{threshold}_{\text{amr}}$, $\gamma_{\text{hydro}}$, $\gamma_{\text{mag}}$, $\kappa$). The variational analysis of the gradient yielded two main results: the demonstration that optimisation requires the action of the mixer and, above all, explicit expressions for the gradient (Equation \ref{eq:gradient-force}) that neatly link the physical quantities of the plasma to the direction and amplitude of the qubit rotations.\\

The Q-HAS framework was implemented as a complete software package, validated through convergence tests ($err < 10^{-10}$ agreement with DNS at maximum depth), six end-to-end QAOA pipeline tests, eight isolated anomaly study cases, and cross-anomaly isolation tests. Eight physical scenarios are supported.

The computational evaluation yields the following quantitative conclusions.  Q-HAS achieves a 0.66\% advantage in composite training loss (0.2134 vs 0.2148) over 170+ Optuna trials.  Captured-fraction advantages are statistically significant on Kelvin--Helmholtz ($\Delta = +0.004$, $p = 0.026$, $d = +0.65$) and Harris Tearing ($\Delta = +0.003$, $p = 0.001$, $d = +2.48$), but also statistically significant \textit{disadvantages} on Rotor ($\Delta = -0.010$, $d = -0.53$) and OT ($\Delta = -0.096$, $d = -1.60$).  On per-cell decision accuracy, classical AMR is correct on 46.6\% of cells vs 43.3\% for Q-HAS (996 vs 925 out of 2136 cells; Fig.~\ref{fig:flip-analysis}).  Q-HAS produces more fragmented patches than classical AMR on Rotor (compactness $2.1\times$ worse) and OT (component density $2.2\times$ worse), contrary to the theoretical prediction that ZZ coupling would enforce spatial coherence (Fig.~\ref{fig:spatial-coherence}).  Classical AMR achieves lower L2 error vs DNS on all four scenarios (Fig.~\ref{fig:physical-fidelity}), though all errors remain below $10^{-3}$.

The strongest evidence for the Hamiltonian's physics-detection capability is localised.  In topology-rich sub-regions (Fig.~\ref{fig:topological}), Q-HAS achieves $+3.8$pp advantage on Tearing and $+5.3$pp on OT.  ZZ-driven flips are 76.5\% accurate on Tearing (where gradient coupling dominates the Hamiltonian); ZZZZ-driven flips are 66.3\% accurate on OT (where the plaquette term dominates). This localised advantage might not translate to a global one because most of the domain is topologically smooth at the current resolution, and corrections from the Z-bias term (which dominates on Rotor) are on average less accurate than the classical baseline.

\paragraph{Limitations}
All quantum circuits are executed on a classical simulator. The 0.66\% training advantage is marginal and scenario-dependent. The ground-truth definition (refine if GT error $>$ mean GT error) is a methodological choice that affects all accuracy metrics. The QAOA circuit cost (8 qubits, COBYLA $K = 40$) adds wall-clock overhead not captured in the compute ratio.

\paragraph{Further work.}\phantom\\

\noindent It could be interesting to:
\begin{enumerate}
    \item scale to larger VQA grids ($4\times 4$, $8\times 8$) to test whether more coupled cells amplify the measured localised advantage.
    \item train and evaluate at higher Reynolds numbers ($Re/Rm \gg 800$) where classical indicator scores may become noisier. \item encode temporal information directly in $H_{\text{struct}}$ to test whether sub-threshold detection is achievable.
    \item switch from finite differences to finite volumes for improved conservation at AMR interfaces.
    \item execute on real quantum hardware to assess the impact of gate noise.
\end{enumerate}

\noindent In summary, Q-HAS proposes a new quantum minimisation-based Hamiltonian approach for performing AMR: the composition of the Hamiltonian corresponds to the physics of each scenario (dominated by gradients for Tearing, dominated by plaquettes for OT, dominated by biases for Rotor; Fig.~\ref{fig:qaoa-detail}), and the spatial coupling terms detect a real structure in topologically rich regions.  At the current scale ($2\times 2$ VQA, $Re = 400$), however, these corrections do not provide a net overall benefit, as the classical multi-indicator score is already nearly optimal.

%TC:ignore
\chapter*{Acknowledgements}
\addcontentsline{toc}{chapter}{Acknowledgements}

I would like to thank my incredible supervisor, Dr.\ Brian Appelbe, for his guidance, constructive feedback and support throughout this project. To write this report and code the associated software, Claude AI (from Anthropic) was used in the following ways: (1) to help identify and resolve logic errors in the MHD solver modules by providing specific code snippets for debugging and optimization; and (2) to rephrase portions of the manuscript to improve English conciseness and clarity by providing my original drafts for stylistic editing. In accordance with departmental guidance, I have personally verified all AI-generated outputs for technical accuracy and physical correctness.

\begin{thebibliography}{99}

\bibitem{freidberg2014}
J.~P. Freidberg, \textit{Ideal MHD} (Cambridge University Press, 2014).

\bibitem{goedbloed2004}
J.~P.~H. Goedbloed and S.~Poedts, \textit{Principles of Magnetohydrodynamics} (Cambridge University Press, 2004).

\bibitem{wesson2011}
J.~Wesson, \textit{Tokamaks}, 4th ed. (Oxford University Press, 2011).

\bibitem{chandrasekhar1961}
S.~Chandrasekhar, \textit{Hydrodynamic and Hydromagnetic Stability} (Clarendon Press, Oxford, 1961).

\bibitem{furth1963}
H.~P. Furth, J.~Killeen, and M.~N. Rosenbluth, ``Finite-resistivity instabilities of a sheet pinch,'' \textit{Phys.\ Fluids} \textbf{6}, 459--484 (1963).

\bibitem{bateman1978}
G.~Bateman, \textit{MHD Instabilities} (MIT Press, 1978).

\bibitem{pope2000}
S.~B. Pope, \textit{Turbulent Flows} (Cambridge University Press, 2000).

\bibitem{berger1984}
M.~J. Berger and J.~Oliger, ``Adaptive mesh refinement for hyperbolic partial differential equations,'' \textit{J.\ Comput.\ Phys.} \textbf{53}, 484--512 (1984).

\bibitem{berger1989}
M.~J. Berger and P.~Colella, ``Local adaptive mesh refinement for shock hydrodynamics,'' \textit{J.\ Comput.\ Phys.} \textbf{82}, 64--84 (1989).

\bibitem{nielsen2010}
M.~A. Nielsen and I.~L. Chuang, \textit{Quantum Computation and Quantum Information}, 10th anniversary ed. (Cambridge University Press, 2010).

\bibitem{harrow2009}
A.~W. Harrow, A.~Hassidim, and S.~Lloyd, ``Quantum algorithm for linear systems of equations,'' \textit{Phys.\ Rev.\ Lett.} \textbf{103}, 150502 (2009).

\bibitem{feynman1982}
R.~P. Feynman, ``Simulating physics with computers,'' \textit{Int.\ J.\ Theor.\ Phys.} \textbf{21}, 467--488 (1982).

\bibitem{farhi2014}
E.~Farhi, J.~Goldstone, and S.~Gutmann, ``A quantum approximate optimization algorithm,'' arXiv:1411.4028 (2014).

\bibitem{preskill2018}
J.~Preskill, ``Quantum computing in the NISQ era and beyond,'' \textit{Quantum} \textbf{2}, 79 (2018).

\bibitem{lucas2014}
A.~Lucas, ``Ising formulations of many NP problems,'' \textit{Front.\ Phys.} \textbf{2}, 5 (2014).

\bibitem{mcclean2018}
J.~R. McClean, S.~Boixo, V.~N. Smelyanskiy, R.~Babbush, and H.~Neven, ``Barren plateaus in quantum neural network training landscapes,'' \textit{Nat.\ Commun.} \textbf{9}, 4812 (2018).

\bibitem{goodfellow2016}
I.~Goodfellow, Y.~Bengio, and A.~Courville, \textit{Deep Learning} (MIT Press, 2016).

\bibitem{powell1994}
M.~J.~D. Powell, ``A direct search optimization method that models the objective and constraint functions by linear interpolation,'' in \textit{Advances in Optimization and Numerical Analysis}, edited by S.~Gomez and J.-P.~Hennart (Springer, 1994), pp.~51--67.

\bibitem{akiba2019}
T.~Akiba, S.~Sano, T.~Yanase, T.~Ohta, and M.~Koyama, ``Optuna: A next-generation hyperparameter optimization framework,'' in \textit{Proc.\ 25th ACM SIGKDD} (2019), pp.~2623--2631.

\bibitem{qiskit2024}
Qiskit contributors, ``Qiskit: An open-source framework for quantum computing,'' \url{https://github.com/Qiskit/qiskit} (2024).

\bibitem{numpy2020}
C.~R. Harris \textit{et al.}, ``Array programming with NumPy,'' \textit{Nature} \textbf{585}, 357--362 (2020).

\bibitem{scipy2020}
P.~Virtanen \textit{et al.}, ``SciPy 1.0: fundamental algorithms for scientific computing in Python,'' \textit{Nat.\ Methods} \textbf{17}, 261--272 (2020).

\bibitem{chorin1968}
A.~J. Chorin, ``Numerical solution of the Navier--Stokes equations,'' \textit{Math.\ Comput.} \textbf{22}, 745--762 (1968).

\bibitem{orszag1979}
S.~A. Orszag and C.~M. Tang, ``Small-scale structure of two-dimensional magnetohydrodynamic turbulence,'' \textit{J.\ Fluid Mech.} \textbf{90}, 129--143 (1979).

%\bibitem{brandt1977}
%A.~Brandt, ``Multi-level adaptive solutions to boundary-value problems,'' \textit{Math.\ Comput.} \textbf{31}, 333--390 (1977).

\bibitem{fryxell2000}
B.~Fryxell \textit{et al.}, ``FLASH: An adaptive mesh hydrodynamics code for modeling astrophysical thermonuclear flashes,'' \textit{Astrophys.\ J.\ Suppl.\ Ser.} \textbf{131}, 273--334 (2000).

\bibitem{stone2008}
J.~M. Stone, T.~A. Gardiner, P.~Teuben, J.~F. Hawley, and J.~B. Simon, ``Athena: a new code for astrophysical MHD,'' \textit{Astrophys.\ J.\ Suppl.\ Ser.} \textbf{178}, 137--177 (2008).

\bibitem{solon2015}
A.~P. Solon \textit{et al.}, ``Pressure is not a state function for generic active fluids,'' \textit{Phys.\ Rev.\ E} \textbf{92}, 062111 (2015).

\bibitem{lohner1987}
R.~L\"ohner, ``An adaptive finite element scheme for transient problems in CFD,'' \textit{Comput.\ Methods Appl.\ Mech.\ Eng.} \textbf{61}, 323--338 (1987).

\bibitem{okubo1970}
A.~Okubo, ``Horizontal dispersion of floatable particles in the vicinity of velocity singularities such as convergences,'' \textit{Deep-Sea Res.} \textbf{17}, 445--454 (1970).

\bibitem{weiss1991}
J.~Weiss, ``The dynamics of enstrophy transfer in two-dimensional hydrodynamics,'' \textit{Physica D} \textbf{48}, 273--294 (1991).

\bibitem{mignone2007}
A.~Mignone, P.~Rossi, G.~Bodo, A.~Ferrari, and S.~Massaglia, ``PLUTO: A numerical code for computational astrophysics,'' \textit{Astrophys.\ J.\ Suppl.\ Ser.} \textbf{170}, 228--242 (2007).

\end{thebibliography}
%TC:endignore

\end{document}